

\documentclass[english, online, cse, master]{dithesis}



\usepackage[type={CC}, modifier={by}, version={4.0}]{doclicense-di}

\IfLang{\MainLang}{finnish}{\copyrighttext
  {\copyrightmetastring\ Jos ei toisin mainita, tämän dokumentin käyttö
    on sallittu Creative Commons 4.0 Kansainvälinen (CC-BY-4.0)
  lisenssin mukaisesti.}
  {\copyrightstring
    \doclicenseThis
    \blockpar
    Joidenkin kuvien ehdot poikkeavat tästä. Poikkeavat ehdot on
    esitetty kuvateksteissä. Näiden kuvien uudelleenjulkaisu tai
    muokkaaminen on sallittu niiden omien ehtojen mukaisesti.
  }
}
{\copyrighttext {\copyrightmetastring\ Except otherwise noted, the reuse of this
    document is authorized under a Creative Commons 4.0 International
  (CC-BY-4.0) license.}{\copyrightstring
    \doclicenseThis
    \blockpar
    Some figures are subject to different terms, such as copyright
    licenses and publisher permissions. There terms are indicated in
    the figure captions. Reproduction or modification of such figures
    is allowed under the conditions specified.
  }
}
 
\title{Toward an Executable ESG Aspect-Based Sentiment Analysis Framework for Indonesian Sustainability Reports}
\otsikko{Kohti toteutettavaa ESG-aspektipohjaisen sentimenttianalyysin viitekehystä Indonesian vastuullisuusraporteille}

\firstname{Daris Dzakwan}
\lastname{Hoesien}



\abstract{This thesis develops an executable framework for analysing Indonesian sustainability reports as record-level evidence rather than as undifferentiated full documents. The study addresses the technical problem of converting heterogeneous, bilingual PDF reports into structured and auditable Environmental, Social, and Governance disclosure records. The implemented workflow combines OCR-based document expansion, page-level provenance tracking, large language model extraction, record-level schema design, ontology alignment, and comparison with ClimateBERT-style climate labels. The framework was evaluated on an active thesis subset of 23 processed reports covering approximately 5,512 pages. The resulting pipeline produced 332 tone-bearing ESG records and 2,074 downstream analytical rows, while all 52 tracked aspects were mapped to ontology paths. The results show that disclosure tone is measurable and useful: commitment was the dominant tone, environmental content was the dominant pillar, and tone labels reached 83.7\% agreement with ClimateBERT-style commitment labels with Cohen's kappa of 0.645. At the same time, the experiments show that extraction quality depends strongly on prompt design and model compliance, with missing-tone cases and schema drift remaining the main weaknesses. The study concludes that tone-aware, provenance-preserving ESG disclosure analysis is feasible and informative, but stronger expert annotation and validation are still required before broader benchmarking claims can be made.}

\tiivistelma{Tässä työssä kehitetään suoritettava kehys indonesialaisten vastuullisuusraporttien analysointiin tietuepohjaisena evidenssinä sen sijaan, että raportteja käsiteltäisiin yhtenäisinä kokonaisuuksina. Työ ratkaisee teknisen ongelman, jossa heterogeeniset ja kaksikieliset PDF-muotoiset raportit muunnetaan rakenteisiksi ja auditoitaviksi Environmental, Social and Governance -sisältöä kuvaaviksi tietueiksi. Toteutettu työnkulku yhdistää OCR-pohjaisen dokumenttien laajennuksen, sivutasoisen alkuperän jäljitettävyyden, suurten kielimallien avulla tehdyn poiminnan, tietuetasoisen skeeman, ontologiaan kohdistamisen sekä vertailun ClimateBERT-tyyppisiin ilmastoluokituksiin. Kehystä arvioitiin työn aktiivisella osajoukolla, joka sisälsi 23 käsiteltyä raporttia ja noin 5 512 sivua. Putki tuotti 332 sävyluokiteltua ESG-tietuetta ja 2 074 jatkoanalyysin riviä, ja kaikki 52 seurattua aspektia pystyttiin liittämään ontologiapolkuihin. Tulokset osoittavat, että ilmaisun sävy on mitattava ja hyödyllinen: sitoumussävy oli yleisin, ympäristö oli yleisin ESG-ulottuvuus, ja sävyluokat saavuttivat 83,7 prosentin yhtäpitävyyden ClimateBERT-tyyppisten sitoumusluokkien kanssa Cohenin kappan ollessa 0,645. Samalla kokeet osoittavat, että poiminnan laatu riippuu vahvasti kehotteiden suunnittelusta ja mallin skeemanoudatuksesta, ja suurimmat heikkoudet liittyvät puuttuviin sävyluokkiin ja skeemavirheisiin. Työn johtopäätös on, että sävytietoinen ja jäljitettävyyden säilyttävä ESG-analyysi on toteuttamiskelpoinen ja informatiivinen, mutta vahvempi asiantuntija-annotointi ja validointi ovat edelleen tarpeen ennen laajempia vertailuväitteitä.}

\keywords{ESG disclosure analysis\spc sustainability reports\spc OCR pipeline\spc tone classification\spc ontology alignment}
\avainsanat{yritysraporttien analytiikka; OCR-putki; sävyluokittelu; ontologiakohdistus; viherpesun seulonta}


\headertitlecap

\addbibresource{citations.bib}

\providecommand{\tightlist}{\setlength{\itemsep}{0pt}\setlength{\parskip}{0pt}}

\begin{document}
\pagestyle{empty}
\maketitlepage
\makecopyrightpage

\IfLang{\MainLang}{finnish}{
  \avainsanat{diplomi-insinöörin tutkinto\spc diplomityön ohjeet\spc diplomityön rakenne\footnote{Käytä avainsanoina eri sanoja kuin lopputyön otsikossa. Poista tämä alaviite lopputyöstäsi.}}
\begin{thesisabstract}{finnish}
\tiivistelmametadata
\end{thesisabstract}
   \keywords{corporate disclosure analytics; OCR pipeline; tone classification; ontology mapping; greenwashing screening}

\begin{thesisabstract}{english}
  \abstractmetadata
\end{thesisabstract}
   \selectlanguage{finnish}
} {
  \keywords{corporate disclosure analytics; OCR pipeline; tone classification; ontology mapping; greenwashing screening}

\begin{thesisabstract}{english}
  \abstractmetadata
\end{thesisabstract}
   \avainsanat{diplomi-insinöörin tutkinto\spc diplomityön ohjeet\spc diplomityön rakenne\footnote{Käytä avainsanoina eri sanoja kuin lopputyön otsikossa. Poista tämä alaviite lopputyöstäsi.}}
\begin{thesisabstract}{finnish}
\tiivistelmametadata
\end{thesisabstract}
   \selectlanguage{english}
}

\contents \thispagestyle{empty}
\header{\headerforeword}

The exploration of Applied NLP in Aspect-based Sentiment Analysis have been an interesting exploration towards building scalable systems with research-grounding exploration. Not only the journey taught me in research perspective, but building real dashboard and application toward the contribution of ESG (Environmental, Social and Governance) field in the perspective of NLP

I owe tremendous gratitude to my supervisors, Dr. Mourad Oussalah and Mehrdad Rostami at the Center for Machine Vision and Signal Analysis (CMVS) at the University of Oulu. Your unwavering support, guidance and expertise provided the perfect foundation for this applied research work, while the freedom to explore unconventional approaches, combined with rigorous scientific standards,
fundamentally shaped my research methodology

I am profoundly grateful to my family, whose love, support, and unwavering belief
in my abilities have been my constant source of strength. Your encouragement
during the most challenging phases of this research, your understanding during long
hours of work, and your celebration of every small milestone have made this journey
meaningful. You have been my anchor throughout this academic voyage.

I gratefully acknowledge the financial support from the University of Oulu Scholarship. This support provided essential
computational resources and the freedom to pursue innovative research directions.

This thesis integrate the research-exploration with the strong curiosity-driven research, and perseverance in facing complex challenges. I hope this work opens exploration on applied solution with research-based that aided the current high demand on ESG implementation to the company 
\signature
 \header{\headerabbreviations}
\setlongtables
\begin{longtable}[l]{p{3cm}p{0.7\textwidth}}

ABSA & aspect-based sentiment analysis\\
    API & application programming interface\\
    CER & character error rate\\
    ESG & environmental, social, and governance\\
    GUI & graphical user interface \\
    JSON & JavaScript Object Notation \\
    JSONL & JSON Lines \\
    LLM & large language model \\
    NLP & natural language processing \\
    OCR & optical character recognition \\
    PDF & portable document format  \\    
    TF-IDF & term frequency-inverse document frequency \\
    WER & word error rate
\end{longtable}
\setcounter{table}{0}

     
\mainmatter
\listoffigures
\listoftables

\chapter{Introduction} \label{introduction} 
\section{Motivation}

Environmental, Social, and Governance (ESG) disclosure has transitioned from a niche voluntary practice to a critical requirement for investors, regulators, and researchers seeking to evaluate corporate sustainability performance  \parencite{hans_bonde_christensen_45b44695}. However, the current landscape of sustainability reporting is characterized by extreme heterogeneity. Reports vary significantly in format and extent due to a lack of global standardization, often blending concrete performance data with vague, forward-looking narratives  \parencite{habiba_al_shaer_c2f5394b}.

In the Indonesian context, these challenges are further compounded. Listed companies often produce bilingual reports with varying degrees of transparency, and empirical evidence suggests that the readability of Indonesian sustainability reports remains low  \parencite{desi_adhariani_543335d6}. This lack of clarity often obscures specific sustainability issues, creating a gap between disclosed intentions and actual corporate conduct  \parencite{marsdenia_010e13e7}.

Traditional document-level analysis—which assigns a single score or sentiment to an entire report—is increasingly viewed as insufficient. ESG content is inherently multifaceted; a single page might contain an environmental commitment, a social program update, and a generic governance compliance statement  \parencite{keane_ong_0a5de439}. This complexity allows for "greenwashing," where companies use "soft language" to mask a lack of measurable outcomes  \parencite{keane_ong_8fec4760}. To address this, this research argues that ESG disclosures must be treated as a record-level \textbf{Aspect-Based Sentiment Analysis} problem. By treating every individual statement as a discrete record classified by aspect, pillar, sentiment, and provenance, we can transform raw PDF corpora into aligned, inspectable, and evaluable computational artifacts.
\section{Problem Statement}

The central research problem is the absence of a reproducible and validated pipeline for transforming heterogeneous sustainability-report PDFs into structured, record-level evidence. Current manual review processes are unscalable, while standard automated tools often fail to capture the semantic nuance required for ESG analysis.

This research addresses six critical technical challenges:
\begin{enumerate}
\item 

\textbf{Extraction Provenance:} Maintaining a clear link between extracted text and its original page/document location in heterogeneous PDFs.\item 

\textbf{Prompt Configuration:} Designing Large Language Model prompts that can reliably extract multi-faceted ESG statements across different providers.\item 

\textbf{Granular Labeling:} Assigning precise aspect, pillar, and tone labels at the record level rather than the document level.\item 

\textbf{Construct Validity:} Identifying the divergence between general sentiment/tone and domain-specific climate labels (e.g., ClimateBERT).\item 

\textbf{Quality Auditing:} Systematically identifying failures in OCR, schema drift, and missing data fields.\item 

\textbf{Evidence Visualization:} Converting large-scale extraction results into interactive dashboards and semantic graphs for research verification.
\end{enumerate}
\section{Research Questions}

The study is guided by the following research questions:
\begin{itemize}
\item 

\textbf{RQ1. ESG ABSA Schema:} How can ESG disclosures be effectively represented using a record-level schema that integrates aspect, ESG pillar, sentiment, and tone?\item 

\textbf{RQ2. Tone vs. Climate-Specific Models:} How do LLM-generated tone labels compare with specialized outputs from models like ClimateBERT, and what does this reveal about the validity of automated ESG assessments?\item 

\textbf{RQ3. Pipeline Diagnostics:} What specific failure modes, such as OCR-related data loss or ontology gaps, characterize the automated extraction of ESG records?\item 

\textbf{RQ4. Stability and Reproducibility:} To what extent do ESG extraction outputs remain stable across varying prompts, LLM models, and service providers?
\end{itemize}

\section{Revision Focus}

This revised manuscript addresses four issues that are critical to the thesis contribution. First, the methodology now states more explicitly that quantitative ESG evidence cannot be limited to narrative text alone, and therefore requires a dedicated table-preservation path during ingestion and preprocessing. Second, the evaluation framing is tightened to acknowledge that Indonesian-language performance cannot be inferred from multilingual models alone, which motivates an Indonesian-specific baseline in the experimental design. Third, the results chapter is positioned around completed evidence only, with unsupported placeholders and unfinished presentation elements removed from the main argument. Fourth, the discussion narrows the theoretical claim of the tone layer: it is presented as a record-level disclosure-stage distinction that complements, rather than replaces, existing disclosure-quality and sustainability-intention frameworks.

These revisions do not change the overall objective of the thesis. They clarify its evidential boundaries, strengthen the methodological justification for the Indonesian ESG context, and align the written claims more closely with what the current executable pipeline can substantiate.

\section{Objectives and Contributions}
\subsection{Research Objectives:}
\begin{itemize}
\item 

\textbf{O1: To Build a Reproducible OCR-to-ESG Extraction Pipeline:} Develop a system that converts raw PDF reports into usable text while preserving page-level provenance for auditability.\item 

\textbf{O2: To Design a Record-Level ESG ABSA Schema:} Establish a technical framework for classifying extracted statements based on granular ESG pillars, sentiment dimensions, and disclosure-stage tone distinctions.\item 

\textbf{O3: To Evaluate Model Divergence:} Conduct a comparative analysis between general LLM tone labels, Indonesian-oriented baselines, and specialized ClimateBERT climate labels to assess construct validity.\item 

\textbf{O4: To Create a Diagnostic Framework:} Implement automated audits for parsing failures, schema drift, and missing labels within the extraction pipeline.\item 

\textbf{O5: To Develop a Semantic Evidence Layer:} Map extracted aspects to a formal ESG ontology and provide interactive dashboards for real-time data exploration and evidence verification.
\end{itemize}
\subsection{Expected Contributions:}
\begin{itemize}
\item 

\textbf{Technical Framework:} An executable, modular pipeline for automated ESG record extraction from sustainability reports.\item 

\textbf{Annotation Schema:} A bilingual, record-level ESG schema tailored for the Indonesian reporting context.\item 

\textbf{Evidence Layer:} A linked data structure connecting raw reports to LLM metadata, human annotations, and validation tables.\item 

\textbf{Methodological Insights:} A comprehensive stability analysis of LLM prompts, a clearer treatment of quantitative table evidence, and a critical discussion of why disclosure-stage tone and climate-specific commitment labels are related but non-interchangeable.
\end{itemize}
\section{Thesis Structure}
\begin{itemize}
\item 

\textbf{Chapter 1:} Outlines the motivation, problem statement, and technical objectives of the study.\item 

\textbf{Chapter 2:} Examines the evolution of NLP in sustainability, the challenge of greenwashing, and the current state-of-the-art in LLM-based extraction.\item 

\textbf{Chapter 3:} Details the OCR-to-ABSA pipeline, including prompt engineering strategies, table-preserving ingestion logic, and the validation framework using ClimateBERT and human annotation.\item 

\textbf{Chapter 4:} Presents implementation details, evaluation metrics, Indonesian-baseline framing, and comparative results regarding model stability and failure reduction.\item 

\textbf{Chapter 5:} Evaluates the ESG ABSA schema, analyzes failure modes, and discusses the implications and boundaries of the disclosure-stage tone framework for future ESG research.\item 

\textbf{Chapter 6:} Summarizes the findings and suggests future directions for automated sustainability analysis.
\end{itemize}
 

\chapter{Related Works}
  \label{relatedwork}
  \section{ Corporate Sustainability Reporting and the Greenwashing Challenge}
\subsection{ The Technical Complexity of ESG Disclosures}

Sustainability reports are inherently difficult to process because they are typically semi-structured PDF documents containing a mix of qualitative narratives, quantitative tables, and visual imagery. Technically, the challenge begins with \textbf{high-fidelity data extraction}; many reports utilize complex layouts that require robust OCR and parsing pipelines to ensure that textual data is not lost or corrupted during the ingestion phase  \parencite[]{marco_bronzini_bf94d949, jacob_beck_9548a833}. Research indicates that data gaps in reporting often stem from these extraction failures rather than a lack of information, making the preprocessing pipeline a critical precursor to greenwashing detection  \parencite{jacob_beck_9548a833}.
\subsection{ NLP Frameworks for Greenwashing Detection}

Recent literature has moved toward a \textbf{multi-dimensional textual framework} to identify greenwashing. This involves analyzing specific technical features:
\begin{itemize}
\item 

\textbf{Thematic and Sentiment Features:} Studies have shown that greenwashing is often characterized by a high frequency of "positive-sentiment" words paired with a lack of concrete, thematic "action" words  \parencite{xiaojia_wang_02351a60}. High-risk reports typically exhibit a "sentimental mask," where optimistic language is used to hide a lack of specific environmental targets  \parencite{xiaojia_wang_02351a60}.\item 

\textbf{"Cheap Talk" and Cherry-Picking:} Advanced models like \textbf{ClimateBERT} have been used to identify "cheap talk"—vague, non-committal disclosures—and "cherry-picking," where firms only report on positive climate-related risks while omitting material negative ones  \parencite{julia_bingler_268445e3}. Technically, this is achieved by comparing a firm's disclosure patterns against industry-wide benchmarks or historical data to find significant deviations  \parencite{julia_bingler_268445e3}.
\end{itemize}
\subsection{ Aspect-Action Analysis and Semantic Inconsistency}

A significant technical advancement in this field is \textbf{Aspect-Action Analysis}. This method goes beyond simple sentiment analysis by extracting specific "aspects" (e.g., carbon emissions reduction) and linking them to "actions" (e.g., "invested \$50M in carbon capture")  \parencite{keane_ong_369a0b63}. Greenwashing is technically identified when there is a \textbf{semantic inconsistency} or a cross-category generalization—such as when a company makes broad environmental claims based on a single, minor action  \parencite{keane_ong_369a0b63}. This requires models with high "cross-category generalization" capabilities to detect if a company's "green" claims in one sector are being used to offset "brown" activities in another  \parencite{keane_ong_369a0b63}.
\subsection{ Monitoring Systems and Empirical Testing}

Newer systems, such as \textbf{DeepGreen}, utilize Large Language Models to create monitoring systems designed for empirical testing  \parencite{chong_xu_bff21576}. These systems are technically designed to:
\begin{enumerate}
\item 

\textbf{Extract structured evidence} from unstructured reports to verify if specific ESG metrics (like Scope 1 or 2 emissions) are reported with sufficient detail  \parencite{chong_xu_bff21576}.\item 

\textbf{Compare internal reporting with external news}, using LLMs to check for discrepancies between what a company says in its ESG report and what is reported in the media  \parencite[]{mariya_pavlova_259be499, fabian_billert_a5fbb687}\item 

\textbf{Score "Greenwashing Tendency"} by applying weighted metrics to the extracted data, identifying when a firm’s rhetoric significantly outpaces its data-backed performance  \parencite{chong_xu_bff21576}.
\end{enumerate}
\subsection{Technical Limitations: The "Soft Language" Problem}

A major hurdle remains the "soft" nature of ESG language. LLMs must be sensitive enough to distinguish between \textbf{substantive disclosures} and \textbf{boilerplate language} that satisfies regulatory checklists without providing material insight  \parencite{frederik_maibaum_b8629037}. This provides the technical justification for \textbf{Section 3.4.4}, as standard machine learning baselines often fail to capture the subtle linguistic nuances that separate genuine commitment from strategic ambiguity  \parencite[]{marco_bronzini_bf94d949, arash_hajikhani_cbb2528f}
\subsection{Table: State-of-the-Art in NLP-Driven ESG and Greenwashing Analysis}

Below is a comprehensive table of state-of-the-art research relevant to Chapter 2.1, focusing on the technical and NLP methodologies used to address greenwashing and ESG extraction. This synthesis highlights how current frameworks transition from simple keyword counting toward complex, intent-aware architectures that verify sustainability claims against verifiable operational evidence  \parencite{keane_ong_0a5de439}. On top of that, there next table focus on the technical nuances of goal alignment, model sensitivity, and sector-specific extraction, complementing the previous literature by addressing the "soft language" and benchmarking challenges in ESG.



\begin{landscape}
\begin{table}[p]
\centering
\caption{State-of-the-Art in the Evolution of NLP for Sustainability}
\label{tab:nlp-evolution-sustainability-1}
\small
\renewcommand{\arraystretch}{1.2}
\begin{tabularx}{\linewidth}{|p{1.8cm}|X|X|p{1.8cm}|X|X|}
\hline
Authors & Method / Approach & Datasets & Modalities Involved & Key Contributions & Limitations \\ \hline
\textbf{Xu et al.}  \parencite{chong_xu_bff21576} & \textbf{DeepGreen System}: LLM-driven monitoring for empirical testing of greenwashing. & Corporate ESG reports and external news articles. & Text & Developed a monitoring system to cross-verify internal reports with external media for discrepancies. & Potential focus on regional reporting standards may limit global generalization. \\ \hline
\textbf{Bingler et al.}  \parencite{julia_bingler_268445e3} & \textbf{ClimateBERT}: Transformer-based model fine-tuned on climate-related data. & Large-scale corporate climate risk disclosures. & Text & Identified "cheap talk" and "cherry-picking" in risk disclosures by detecting vague, non-committal language. & Primarily focused on climate/risk; may overlook social (S) and governance (G) specific greenwashing. \\ \hline
\textbf{Ong et al.}  \parencite{keane_ong_369a0b63} & \textbf{Aspect-Action Analysis}: Cross-category generalization for robust ESG analysis. & ESG reports categorized by specific environmental actions. & Text & Links specific "aspects" to concrete "actions" to identify inconsistencies between rhetoric and investment. & Requires high-quality labeled data for different action categories to maintain accuracy. \\ \hline
\textbf{Wang et al.}  \parencite{xiaojia_wang_02351a60} & \textbf{Thematic \& Sentiment Analysis}: Predicting greenwashing degrees via linguistic features. & Corporate ESG reports. & Text (Sentiment/ Theme) & Quantified how high positive sentiment correlates with a lack of specific, data-backed environmental themes. & Purely linguistic focus may miss technical data discrepancies (e.g., mismatched emission numbers). \\ \hline
\textbf{Bronzini et al.}  \parencite{marco_bronzini_bf94d949} & \textbf{LLM-Structured Extraction}: Deriving structured insights from sustainability reports. & Diverse set of corporate sustainability reports. & Text \& Semi-structured & Demonstrated the superiority of LLMs in extracting structured data from unstructured narrative text. & High computational cost and "hallucination" risks in extracting specific numerical metrics. \\ \hline
\textbf{Beck et al.}  \parencite{jacob_beck_9548a833} & \textbf{Benchmark Dataset}: GHG emission extraction and data gap analysis. & Global dataset of PDF reports for Greenhouse Gas monitoring. & Text \& PDF Metadata & Created a specialized benchmark for extracting Scope 1, 2, and 3 emissions to address reporting data gaps. & Focus is narrow, not covering broader qualitative greenwashing narratives. \\ \hline
\textbf{Schimanski et al.}  \parencite{tobias_schimanski_3549ff9f} & \textbf{Nature-Related Analysis}: Datasets and models for biodiversity and nature disclosures. & Nature-related corporate disclosures. & Text & Expanded the scope of NLP moving beyond simple climate metrics. & Nature-related reporting is still in its infancy, leading to sparse and inconsistent ground-truth data. \\ \hline
\end{tabularx}
\end{table}
\end{landscape}



\begin{landscape}
\begin{table}[p]
\centering
\caption{NLP-Driven ESG and Greenwashing Analysis}
\label{tab:nlp-driven-esg-greenwashing}
\small
\renewcommand{\arraystretch}{1.2}
\begin{tabularx}{\linewidth}{|p{1.8cm}|X|X|p{1.8cm}|X|X|}
\hline
Authors & Method / Approach & Datasets & Modalities Involved & Key Contributions & Limitations \\ \hline
\textbf{Kılınç et al.}  \parencite{yavuz_k_l_n__64e32f80} & \textbf{Multi-dimensional Textual Framework}: Specifically designed for greenwashing detection. & Multi-sector corporate sustainability reports. & Textual narratives & Developed a framework that integrates linguistic indicators of deception with sustainability metrics. & Framework performance is highly dependent on the quality of the underlying sentiment and thematic lexicons. \\ \hline
\textbf{Schimanski et al.}  \parencite{tobias_schimanski_241ada20} & \textbf{ClimateBERT-NetZero}: Specialized model for detecting and assessing Net Zero targets. & Corporate Net Zero pledge documents and reports. & Text (Targets/ Claims) & Created a pipeline to distinguish between "vague" commitments and "verifiable" reduction targets, essential for flagging target-based greenwashing. & Difficulty in handling "future-looking" statements that lack immediate material evidence. \\ \hline
\textbf{Maibaum et al.}  \parencite{frederik_maibaum_b8629037} & \textbf{Tool Selection Framework}: Systematic classification of sustainability information. & Standardized corporate reporting datasets. & Text & Established a decision-matrix for choosing between classical NLP and Deep Learning based on data complexity. & Identifies that traditional tools often fail at capturing the semantic nuance required for complex ESG themes. \\ \hline
\textbf{Chung \& Latifi}  \parencite{tae_sun_chung_e16a4c18} & \textbf{ESG Domain-Specific Evaluation}: Benchmarking domain-specific LLMs against traditional ML. & ESG-specific text classification benchmarks. & Text & Demonstrated that domain-pretraining (e.g., ESG-BERT) significantly outperforms general models in identifying material ESG risks. & Domain-specific models can exhibit high sensitivity to the specific training corpus, leading to potential bias. \\ \hline
\textbf{Yu et al.}  \parencite{yinan_yu_e282fd98} & \textbf{climateBUG}: Data-driven framework for analyzing bank reporting through a climate lens. & Financial sector climate risk reports. & Text \& Financial Data & Developed a sector-specific lens for climate reporting, highlighting the "transparency gap" in banking disclosures. & Sector-specific focus (banking) may not directly transfer to heavy-industry or retail ESG reporting. \\ \hline
\textbf{Hajikhani \& Cole}  \parencite{arash_hajikhani_cbb2528f} & \textbf{Sensitivity and Bias Review}: Critical analysis of LLMs in specialized AI domains. & General and specialized LLM outputs. & Textual outputs & Highlighted the risks of "hallucination" and bias when LLMs are used for specialized tasks like ESG auditing. & Focuses on the "path toward specialized AI" rather than providing a standalone extraction tool. \\ \hline
\end{tabularx}
\end{table}
\end{landscape}
\section{ Evolution of NLP in Sustainability Analysis}

The methodology for analyzing corporate sustainability has undergone a significant transformation, evolving from simple keyword counting to sophisticated semantic understanding. This progression reflects the increasing complexity of ESG disclosures and the need for tools that can detect nuanced linguistic patterns like greenwashing.
\subsection{Dictionary-Based and Statistical Foundations}

Early NLP applications in sustainability relied heavily on \textbf{dictionary-based techniques} and \textbf{lexicons}. These methods used predefined lists of "green" or "socially responsible" keywords to quantify the volume of ESG communication  \parencite{frederik_maibaum_b8629037}. While foundational, research has shown that dictionaries exhibit a large variation in quality and fail to capture the context in which words are used  \parencite{frederik_maibaum_b8629037}. For instance, a dictionary might flag the word "carbon" as a positive environmental indicator, even if the text discusses a "failure to reduce carbon emissions"  \parencite{julia_bingler_268445e3}. Subsequent statistical methods, such as \textbf{Latent Dirichlet Allocation} for topic modeling, offered improved performance over dictionaries by identifying latent themes across reports, yet they remained limited by their inability to handle complex semantic relationships  \parencite{frederik_maibaum_b8629037}.
\subsection{The Transformer Revolution and Contextual Embeddings}

The introduction of the \textbf{Transformer architecture} and models like \textbf{BERT} marked a paradigm shift. Unlike previous word embedding techniques (e.g., Word2Vec) that assigned a single vector to a word regardless of context, Transformers use attention mechanisms to derive context-dependent meanings  \parencite{frederik_maibaum_b8629037}. This capability is crucial for ESG analysis, where the same term can have vastly different implications depending on the surrounding text  \parencite{natraj_raman_009776f2}. Studies demonstrate that fine-tuned Transformer models significantly outperform classical machine learning baselines in classifying material ESG risks and identifying sustainability trends in corporate earnings calls  \parencite[]{tae_sun_chung_e16a4c18, natraj_raman_009776f2}.
\subsection{ Domain-Specific Specialization: The Rise of ClimateBERT}

As the limitations of general-purpose language models became apparent, the field shifted toward \textbf{domain-specific fine-tuning}. A prominent example is \textbf{ClimateBERT}, a model pretrained on millions of climate-related sentences from corporate reports and news  \parencite[]{tobias_schimanski_4e192ec5, julia_bingler_268445e3}.
\begin{itemize}
\item 

\textbf{Precision in "Cheap Talk" Detection:} ClimateBERT has proven highly effective at identifying "cheap talk"—vague, non-committal climate disclosures that satisfy regulatory requirements without committing to material action  \parencite{julia_bingler_268445e3}.\item 

\textbf{Enhanced Performance:} By training on domain-specific corpora, these models achieve much higher accuracy in detecting climate-related risks compared to general models like BERT-base  \parencite[]{tae_sun_chung_e16a4c18, julia_bingler_268445e3}.This specialization has extended to specific sub-domains, with researchers developing distinct models for the "E," "S," and "G" components to capture the unique linguistic styles of each pillar  \parencite{tobias_schimanski_4e192ec5}.
\end{itemize}
\subsection{ Large Language Models and Generative Extraction}

The current frontier in the evolution of ESG NLP is the adoption of \textbf{Large Language Models} such as GPT-4 and Llama. These models represent a move from simple text classification to \textbf{structured semantic extraction}  \parencite{seyed_alireza_mousavian_anaraki_b56a1c0f}.
\begin{itemize}
\item 

\textbf{One-Shot and Few-Shot Capabilities:} Recent evaluations show that LLMs, when paired with well-designed prompts, can surpass earlier fine-tuned models in extracting structured data from unstructured reports  \parencite{frederik_maibaum_b8629037}.\item 

\textbf{Multi-Dimensional Analysis:} LLMs enable more complex tasks, such as \textbf{Aspect-Action Analysis}, which links specific environmental aspects to concrete corporate actions to identify inconsistencies  \parencite{keane_ong_369a0b63}.\item 

\textbf{Expansion to Nature and Biodiversity:} New datasets and models are now emerging to tackle "nature-related" disclosures, moving beyond carbon-centric analysis to evaluate how companies report on biodiversity and ecosystem services  \parencite{tobias_schimanski_3549ff9f}.
\end{itemize}
\subsection{Table: State-of-the-Art in the Evolution of NLP for Sustainability}

The table below summarizes the key milestones in the technical evolution of NLP within sustainability analysis, tracking the shift from lexicon-based systems to domain-specific Transformers and eventually to generative Large Language Models, also the table introduced specialized frameworks and models that represent the transition from general climate classification to more granular, sector-specific, and target-oriented analysis

\begin{landscape}
\begin{table}[p]
\centering
\caption{State-of-the-Art in the Evolution of NLP for Sustainability}
\label{tab:nlp-evolution-sustainability-2}
\small
\renewcommand{\arraystretch}{1.2}
\begin{tabularx}{\linewidth}{|p{1.8cm}|p{3.5cm}|p{3cm}|p{1.8cm}|X|X|}
\hline
Authors & Method / Approach & Datasets & Modalities Involved & Key Contributions & Limitations \\ \hline
\textbf{Maibaum et al.}  \parencite{frederik_maibaum_b8629037} & \textbf{Tool Selection Framework}: Comparing Dictionaries, BERT, and LLMs. & Standardized corporate reporting datasets. & Text & Established a decision-matrix showing that while dictionaries provide transparency, Transformers are required for semantic nuance. & Does not propose a new architecture, but rather a methodology for tool selection. \\ \hline
\textbf{Raman et al.}  \parencite{natraj_raman_009776f2} & \textbf{Distant Supervision}: Neural Language Models for trend mapping. & Corporate earnings call transcripts and ESG news. & Text & One of the first applications of neural embeddings to map long-term ESG trends via distant supervision. & Reliance on weak labels (distant supervision) can introduce noise into trend analysis. \\ \hline
\textbf{Bingler et al.}  \parencite{julia_bingler_268445e3} & \textbf{ClimateBERT}: Domain-specific pre-training. & 1.6M paragraphs of climate-related disclosures. & Text & Proven that general-purpose models fail to detect "cheap talk" and that domain-pretraining is essential for climate precision. & Narrow focus on climate-related text; less effective for social or governance specificities. \\ \hline
\textbf{Schimanski et al.}  \parencite{tobias_schimanski_4e192ec5} & \textbf{BERT-Based Measurement}: Quantifying multi-dimensional ESG communication. & Large-scale sustainability reports across multiple industries. & Text & Bridged the gap between manual content analysis and automated measurement using contextual embeddings. & The model is a "black box" compared to dictionary methods, making it harder for auditors to interpret specific flags. \\ \hline
\textbf{Chung \& Latifi}  \parencite{tae_sun_chung_e16a4c18} & \textbf{Domain-Specific LLM Evaluation}: Benchmarking ESG-specific LLMs. & Specialized ESG text classification benchmarks. & Text & Quantified the performance leap when using models pre-trained specifically on financial and ESG corpora. & Primarily addresses classification tasks rather than complex structured data extraction. \\ \hline
\textbf{Bronzini et al.}  \parencite{marco_bronzini_bf94d949} & \textbf{Generative Extraction}: Transition from classification to structured insights via LLMs. & 100+ high-complexity corporate sustainability reports. & Text \& Semi-structured data & Demonstrated that LLMs (e.g., GPT-4) can derive structured tables and insights that previously required manual entry. & High computational cost and the inherent risk of numerical "hallucinations" in extracted data. \\ \hline
\textbf{Anaraki et al.}  \parencite{seyed_alireza_mousavian_anaraki_b56a1c0f} & \textbf{Systematic Review}: Mapping the research agenda for LLMs in sustainability. & Literature corpus on LLM applications in ESG. & Meta-analysis & Provided a comprehensive taxonomy of how LLMs are currently being used to automate regulatory reporting (GRI/SASB). & As a review paper, it does not provide an empirical benchmark for specific extraction tasks. \\ \hline
\end{tabularx}
\end{table}
\end{landscape}



\begin{landscape}
\begin{table}[p]
\centering
\small
\renewcommand{\arraystretch}{1.2}
\begin{tabularx}{\linewidth}{|p{1.8cm}|p{3.5cm}|p{3cm}|p{1.8cm}|X|X|}
\hline
Authors & Method / Approach & Datasets & Modalities Involved & Key Contributions & Limitations \\ \hline
\textbf{Ong et al.}  \parencite{keane_ong_369a0b63} & \textbf{Aspect-Action Analysis}: Cross-category generalization for robust analysis. & ESG reports with labeled environmental actions. & Text & Developed a technical logic to link broad environmental claims ("Aspects") to verifiable investments ("Actions") to detect narrative inconsistencies  \parencite{keane_ong_369a0b63}. & High performance dependency on the availability of explicit action-oriented language in reports. \\ \hline
\textbf{Schimanski et al.}  \parencite{tobias_schimanski_3549ff9f} & \textbf{Nature-Related NLP}: Domain-specific datasets/models for biodiversity. & TNFD-aligned corporate nature disclosures. & Text & Technically expanded the evolution of ESG NLP from carbon-centric to nature/biodiversity-centric  \parencite{tobias_schimanski_3549ff9f}. & Biodiversity reporting is currently sparse, leading to "small data" challenges for model training. \\ \hline
\textbf{Billert \& Conrad}  \parencite{fabian_billert_a5fbb687} & \textbf{Nano-ESG Extraction}: Information extraction from external news. & Large corpus of corporate news articles. & Text & Shifted from "Internal-Only" to "External-Monitoring" by extracting ESG performance indicators from media to verify company claims  \parencite{fabian_billert_a5fbb687}. & News data is significantly noisier than formal reports, requiring complex cleaning and deduplication pipelines. \\ \hline
\textbf{Schimanski et al.}  \parencite{tobias_schimanski_241ada20} & \textbf{ClimateBERT-NetZero}: Specialized target assessment pipeline. & Net Zero pledge documents and corporate targets. & Text (Targets/ Claims) & Introduced a methodology to technically differentiate between "vague" pledges and "verifiable" science-based targets  \parencite{tobias_schimanski_241ada20}. & Struggles to validate future-looking claims that lack current-year material evidence or milestones. \\ \hline
\textbf{Wang et al.}  \parencite{xiaojia_wang_02351a60} & \textbf{Multi-Feature Prediction}: Sentiment and thematic feature integration. & Standard corporate ESG reports. & Text & Quantified the "Greenwashing Degree" by using the disparity between positive sentiment and specific thematic density as a technical predictor  \parencite{xiaojia_wang_02351a60}. & Focused on linguistic features; may miss numerical discrepancies found in financial tables. \\ \hline
\textbf{Yu et al.}  \parencite{yinan_yu_e282fd98} & \textbf{climateBUG}: Sector-specific framework for financial institutions. & Bank-specific climate risk and TCFD reports. & Text \& Metadata & Developed a data-driven framework specifically for the banking sector to identify "transparency gaps" in carbon-intensive lending  \parencite{yinan_yu_e282fd98}. & The banking-specific ontology lacks generalization to other sectors like heavy manufacturing or retail. \\ \hline
\end{tabularx}
\end{table}
\end{landscape}



\section{Large Language Models for Structured ESG Extraction}

The shift toward Large Language Models represents a fundamental technical pivot from \textbf{text classification}—where models simply categorize a paragraph as "Environmental"—to \textbf{structured information extraction}, where models extract specific, actionable data points into predefined schemas  \parencite[]{seyed_alireza_mousavian_anaraki_b56a1c0f, marco_bronzini_bf94d949}.
\subsection{From Classification to Schema-Based Extraction}

Traditional NLP models typically perform sequence classification, providing a high-level label for a text block. In contrast, LLMs enable the derivation of structured insights, such as transforming unstructured narrative text into JSON or tabular formats that include specific metrics like Scope 1 emissions, target years, and baseline comparisons  \parencite{marco_bronzini_bf94d949}. This capability is critical for automating the assessment of sustainability reports, which are often composed of semi-structured data that varies significantly between companies  \parencite[]{jacob_beck_9548a833, marco_bronzini_bf94d949}. Research by \textbf{Bronzini et al.} demonstrates that LLMs like GPT-4 can derive these structured tables with a level of precision that previously required manual auditing, though the risk of numerical "hallucinations" remains a technical hurdle  \parencite{marco_bronzini_bf94d949}.
\subsection{Prompt Engineering and In-Context Learning}

A key technical advantage of LLMs in the ESG domain is their ability to perform tasks via \textbf{In-Context Learning}, which includes zero-shot and few-shot prompting strategies  \parencite{frederik_maibaum_b8629037}. This eliminates the need for the extensive, task-specific fine-tuning required by earlier models like BERT  \parencite[]{tae_sun_chung_e16a4c18, frederik_maibaum_b8629037}.
\begin{itemize}
\item 

\textbf{Zero-Shot Extraction:} LLMs can extract ESG data based solely on a natural language description of the schema (aligning with \textbf{Section 3.4.1 Schema Layer})  \parencite{marco_bronzini_bf94d949}.\item 

\textbf{Few-Shot Prompting:} Providing a few expert-labeled examples (e.g., how to identify a "Science-Based Target") within the prompt significantly improves model performance on complex, "soft" language tasks  \parencite[]{frederik_maibaum_b8629037, arash_hajikhani_cbb2528f}.\item 

\textbf{Chain-of-Thought:} Encouraging the model to "reason" through an extraction—such as explaining \textit{why} a specific sentence constitutes a carbon reduction action—helps mitigate errors in judgment and improves the interpretability of the results  \parencite{arash_hajikhani_cbb2528f}.
\end{itemize}
\subsection{Aspect-Action Linkage and Semantic Mapping}

For greenwashing detection, simple keyword matching is insufficient. LLMs excel at \textbf{Aspect-Action Analysis}, a technical logic where the model identifies an environmental "aspect" (e.g., water conservation) and attempts to link it to a verifiable "action" (e.g., \$10M investment in desalination plants)  \parencite{keane_ong_369a0b63}. Technically, this involves complex semantic mapping where the LLM evaluates the strength of the connection between a claim and the evidence provided. If a model identifies a high frequency of "aspects" without corresponding "actions," it can flag the report for potential greenwashing based on narrative inconsistency  \parencite[]{xiaojia_wang_02351a60, keane_ong_369a0b63}.
\subsection{ Domain-Specific vs. General LLMs}

While general LLMs like GPT-4 show high performance, the evolution of the field has led to \textbf{domain-specific pre-trained models} (e.g., ESG-BERT or specialized Llama variants). Studies comparing these models indicate that domain-specific pre-training on financial and corporate corpora allows the model to better understand the "jargon" of sustainability reporting, leading to higher accuracy in identifying material risks  \parencite[]{seyed_alireza_mousavian_anaraki_b56a1c0f, tae_sun_chung_e16a4c18}. However, general-purpose LLMs often retain a lead in "reasoning" tasks and complex instruction following, which justifies the \textbf{Comparative Model Logic} found in \textbf{Section 3.4.5}  \parencite[]{arash_hajikhani_cbb2528f, tae_sun_chung_e16a4c18}.
\subsection{Technical Risks: Hallucinations and Grounding}

A critical limitation of LLMs in ESG extraction is their sensitivity to input phrasing and their tendency to generate plausible-sounding but factually incorrect data—commonly known as \textbf{hallucinations}  \parencite{arash_hajikhani_cbb2528f}. In the context of sustainability, a hallucinated emission figure or target date could lead to incorrect regulatory filings. This technical risk necessitates the use of \textbf{External Validators} (like ClimateBERT, as seen in  \textbf{Section 3.4.3}) and \textbf{Human-in-the-Loop} workflows to verify the extracted "silver labels" before they are accepted as ground truth  \parencite[]{fabian_billert_a5fbb687, marco_bronzini_bf94d949, yinan_yu_e282fd98}.

\begin{landscape}
\begin{table}[p]
\centering
\caption{Hallucinations and Grounding}
\label{tab:hallucinations-grounding}
\small
\renewcommand{\arraystretch}{1.2}
\begin{tabularx}{\linewidth}{|p{1.8cm}|p{3.5cm}|p{3.5cm}|p{1.8cm}|X|X|}
\hline
Authors & Method / Approach & Datasets & Modalities Involved & Key Contributions & Limitations \\ \hline
\textbf{Bronzini et al.}  \parencite{marco_bronzini_bf94d949} & \textbf{Structured Parsing}: Using LLMs to derive JSON/tabular insights from text. & 100+ high-complexity corporate sustainability reports. & Text \& Semi-structured data & Demonstrated that LLMs can automate the conversion of qualitative narratives into structured ESG metrics. & High risk of numerical "hallucinations" and significant computational costs per report. \\ \hline
\textbf{Anaraki et al.}  \parencite{seyed_alireza_mousavian_anaraki_b56a1c0f} & \textbf{Systematic LLM Review}: Mapping the technical landscape of LLM applications in ESG. & Comprehensive corpus of LLM-ESG literature. & Meta-analysis & Provided a taxonomy of LLM tasks in sustainability, including regulatory reporting and audit automation. & Does not provide a specific empirical benchmark or new architecture. \\ \hline
\textbf{Hajikhani \& Cole}  \parencite{arash_hajikhani_cbb2528f} & \textbf{Sensitivity Analysis}: Evaluating model bias and sensitivity in specialized AI. & Multi-domain LLM outputs (general and specialized). & Text & Highlighted the critical need for "Self-Refinement" and specialized prompts to mitigate bias in LLM outputs. & Focuses on the "path toward specialized AI" rather than a standalone ESG extraction tool. \\ \hline
\textbf{Chung \& Latifi}  \parencite{tae_sun_chung_e16a4c18} & \textbf{Benchmarking ESG-LLMs}: Comparing domain-specific LLMs against general models. & ESG-specific text classification and extraction benchmarks. & Text & Quantified the performance gap, showing that ESG-pre-trained LLMs significantly reduce false-positive greenwashing flags. & Domain-specific models can inherit biases from their specific training corpus. \\ \hline
\textbf{Xu et al.}  \parencite{chong_xu_bff21576} & \textbf{DeepGreen System}: LLM-driven monitoring for empirical greenwashing testing. & Corporate ESG reports and external news articles. & Text & Created a cross-verification pipeline using LLMs to check report validity against external evidence. & Regional data focus (e.g., China) may affect generalizability to global reporting standards. \\ \hline
\textbf{Ong et al.}  \parencite{keane_ong_369a0b63} & \textbf{Aspect-Action Analysis}: Linking sustainability "Aspects" to concrete "Actions." & ESG reports with labeled environmental investments. & Text & Developed a semantic logic to detect greenwashing when claims are not backed by verifiable data. & Highly dependent on the presence of explicit, action-oriented language within the reports. \\ \hline
\textbf{Beck et al.}  \parencite{jacob_beck_9548a833} & \textbf{GHG Benchmark}: Structured extraction of Greenhouse Gas emissions. & Global dataset of PDF sustainability reports. & Text \& PDF Metadata & Established a technical benchmark for extracting precise Scope 1, 2, and 3 emission figures from unstructured text. & Focused exclusively on GHG metrics, omitting broader qualitative ESG themes. \\ \hline
\end{tabularx}
\end{table}
\end{landscape}


\begin{landscape}
\begin{table}[p]
\centering
\caption{Advanced benchmarks and specialized extraction settings for ESG and greenwashing analysis}
\label{tab:advanced-esg-greenwashing-benchmarks}
\small
\renewcommand{\arraystretch}{1.2}
\begin{tabularx}{\linewidth}{|p{1.8cm}|p{3.5cm}|p{2.5cm}|p{1.8cm}|X|X|}
\hline
Authors & Method / Approach & Datasets & Modalities Involved & Key Contributions & Limitations \\ \hline
\textbf{Schimanski et al.}  \parencite{tobias_schimanski_241ada20} & \textbf{ClimateBERT-NetZero}: Targeted assessment of reduction pledges. & Corporate Net Zero commitment documents. & Text (Targets/ Pledges) & Developed a pipeline to extract and assess "Net Zero" targets, distinguishing between vague pledges and verifiable, science-based reduction goals  \parencite{tobias_schimanski_241ada20}. & Often struggles to validate purely forward-looking statements that lack current-year material evidence. \\ \hline
\textbf{Bingler et al.}  \parencite{julia_bingler_268445e3} & \textbf{Domain-Specific Risk Extraction}: Identifying "cherry-picking" in disclosures. & TCFD-aligned climate risk reports. & Text & Technically identified "cheap talk" by extracting non-committal language and comparing it against material climate risks  \parencite{julia_bingler_268445e3}. & Narrowly focused on climate-related risks; less effective for social (S) or governance (G) extraction. \\ \hline
\textbf{Schimanski et al.}  \parencite{tobias_schimanski_3549ff9f} & \textbf{Nature-Related Models}: Extracting biodiversity and ecosystem metrics. & TNFD-aligned corporate nature disclosures. & Text (Nature/ Biodiversity) & Expanded the extraction paradigm to nature-related disclosures, providing a framework for monitoring biodiversity impact  \parencite{tobias_schimanski_3549ff9f}. & Nature-related reporting is still emerging, leading to sparse and inconsistent ground-truth data for model training. \\ \hline
\textbf{Billert \& Conrad}  \parencite{fabian_billert_a5fbb687} & \textbf{Nano-ESG Extraction}: Event-based extraction from news media. & Large-scale corpus of corporate news articles. & Text & Shifted from internal reports to "outside-in" monitoring by extracting ESG-relevant events from news to verify internal claims  \parencite{fabian_billert_a5fbb687}. & News data is significantly noisier than formal reports, requiring complex deduplication and filtering. \\ \hline
\textbf{Yu et al.}  \parencite{yinan_yu_e282fd98} & \textbf{climateBUG Framework}: Sector-specific extraction for financial institutions. & Bank-specific climate risk and TCFD reports. & Text \& Metadata & Developed a data-driven framework specifically for the banking sector to identify "transparency gaps" in carbon-intensive lending  \parencite{yinan_yu_e282fd98}. & The sector-specific ontology (banking) may not directly transfer to heavy-industry or manufacturing sectors. \\ \hline
\textbf{Kılınç et al.}  \parencite{yavuz_k_l_n__64e32f80} & \textbf{Multi-dimensional Framework}: Deception detection in textual reporting. & Diverse corporate sustainability reports. & Textual Narratives & technically linked linguistic "deception" cues (e.g., strategic ambiguity) to specific sustainability extraction categories  \parencite{yavuz_k_l_n__64e32f80}. & High performance dependency on the quality and breadth of the underlying sentiment and thematic lexicons. \\ \hline
\textbf{Wang et al.}  \parencite{xiaojia_wang_02351a60} & \textbf{Sentiment-Thematic Integration}: Predicting greenwashing degree via feature disparity. & Standardized corporate ESG reports. & Text & Quantified how high positive sentiment correlates with a lack of specific, data-backed thematic actions, creating a predictive score for greenwashing  \parencite{xiaojia_wang_02351a60}. & Focused primarily on linguistic features; may miss numerical discrepancies found in financial tables. \\ \hline
\end{tabularx}
\end{table}
\end{landscape}

\section{ Data Gaps and Benchmark Design}

The creation of a robust methodology for sustainability analysis is fundamentally constrained by the quality and structure of available data. This section examines the technical "data gaps" identified in current literature and the evolving standards for benchmark design in the ESG domain.
\subsection{Identifying Technical Data Gaps in Sustainability}

Despite the increasing volume of reports, significant gaps prevent automated systems from achieving high accuracy.
\begin{itemize}
\item 

\textbf{Extraction and Parsing Failures:} A primary technical gap is the failure of automated pipelines to ingest complex, semi-structured PDF layouts  \parencite[]{marco_bronzini_bf94d949, jacob_beck_9548a833}. Research by \textbf{Beck et al.} highlights that data gaps in sustainability reporting—particularly regarding greenhouse gas emissions—are often caused by the technical difficulty of extracting data from multi-column tables and non-standardized units within unstructured narratives  \parencite{jacob_beck_9548a833}.\item 

\textbf{The Scope 3 Reporting Gap:} While Scope 1 and 2 emissions are increasingly reported with consistency, Scope 3 reporting remains sparse and technically inconsistent across industries, making it the most difficult metric to benchmark accurately  \parencite{jacob_beck_9548a833}.\item 

\textbf{"Sentimental Masks" and Omissions:} Data gaps are not always accidental; they can be strategic. Literature identifies "cherry-picking" (reporting only positive metrics) as a practice that creates artificial gaps in a firm’s environmental risk profile  \parencite[]{xiaojia_wang_02351a60, julia_bingler_268445e3}. This requires benchmarks that can identify not just what is present, but what has been omitted  \parencite{julia_bingler_268445e3}.
\end{itemize}
\subsection{The Role of Benchmark Datasets in ESG NLP}

A technical benchmark provides a standardized "ground truth" to evaluate model performance. The field is currently shifting from general classification datasets to \textbf{specialized extraction benchmarks}:
\begin{itemize}
\item 

\textbf{Target-Specific Benchmarking:} Newer benchmarks, such as \textbf{ClimateBERT-NetZero}, are designed specifically to evaluate a model's ability to distinguish between vague corporate pledges and verifiable, science-based reduction targets  \parencite{tobias_schimanski_241ada20}.\item 

\textbf{Nature and Biodiversity:} As the focus expands beyond carbon, specialized datasets are emerging for nature-related disclosures, addressing the "small data" challenge where high-quality labeled examples of biodiversity impact are scarce  \parencite{tobias_schimanski_3549ff9f}.\item 

\textbf{News-Based External Validation:} Benchmarks like the \textbf{ESG-FTSE corpus} utilize external news articles to provide an "outside-in" perspective, allowing researchers to check for discrepancies between internal corporate reports and external media narratives  \parencite[]{fabian_billert_a5fbb687, mariya_pavlova_259be499}.
\end{itemize}
\subsection{Technical Requirements for Benchmark Design}

According to recent studies, a defensible ESG benchmark must satisfy three core technical requirements:
\begin{enumerate}
\item 

\textbf{Granular Unit of Analysis:} Benchmarks must transition from report-level summaries to \textbf{paragraph-level extractions} to capture the specific context of claims  \parencite{jacob_beck_9548a833}.\item 

\textbf{Gold Standard Ground Truth:} While "Silver Labels" (generated via weak supervision or LLMs) are useful for initial seeding, a high-fidelity benchmark requires an expert-led \textbf{Human Annotation Workflow} to verify the semantic nuances of sustainability language  \parencite[]{keane_ong_369a0b63, chong_xu_bff21576}.\item 

\textbf{Cross-Category Generalization:} A robust benchmark should test a model's performance across diverse industry sectors (e.g., banking vs. heavy industry) and ESG categories to ensure the methodology is not overfitted to a single domain  \parencite[]{yinan_yu_e282fd98, keane_ong_369a0b63}.
\end{enumerate}

\subsection{Table: State-of-the-Art in ESG Benchmark Design}

For \textbf{Section 2.4: Data Gaps and Benchmark Design}, the following table identifies the most recent state-of-the-art benchmarks and datasets designed to address the "expert-scarcity" and "data-fragmentation" problems in sustainability reporting. These works provide the technical foundation for the \textbf{Benchmark and Ground Truth Design (Section 3.6)}, and also highlights additional state-of-the-art research focusing on the specific technical gaps identified in corporate reporting and the datasets developed to mitigate them. These works provide the scholarly justification to focus on \textbf{Coverage, Stratification, and Resumability (Section 3.6.4)}.


\begin{landscape}
\begin{table}[p]
\centering
\caption{State-of-the-art ESG benchmarks and dataset design patterns supporting the benchmark and ground-truth construction in Section 3.6.}
\label{tab:esg-benchmark-dataset-design}
\small
\renewcommand{\arraystretch}{1.2}
\begin{tabularx}{\linewidth}{|p{1.8cm}|X|X|p{1.8cm}|X|X|}
\hline
Authors & Method / Approach & Datasets & Modalities Involved & Key Contributions & Limitations \\ \hline
\textbf{Beck et al.} \parencite{jacob_beck_9548a833} & \textbf{LLM-powered extraction with human review}: a multi-stage validation pipeline. & 139 sustainability reports from company websites. & Text \& PDF metadata & Created a gold-standard dataset for GHG emissions to validate automated extraction. & Primarily focused on numerical emission metrics; lacks qualitative greenwashing narrative analysis. \\ \hline
\textbf{Sun et al.} \parencite{siqi_sun_f6a51152} & \textbf{ESG-Bench}: QA-based benchmarking for long-context understanding. & Human-annotated QA pairs grounded in real-world reports. & Text & Developed a benchmark for hallucination mitigation in compliance-critical ESG settings. & Complex QA design may limit direct reuse for simpler extraction tasks. \\ \hline
\textbf{Maibaum et al.} \parencite{frederik_maibaum_b8629037} & \textbf{Validity and quality assessment}: comparison of four NLP techniques. & 75,000 manually labelled sentences from 10-K reports. & Text & Established large-scale ground truth showing that dictionaries fail to capture semantic context. & Focuses on 10-K filings, which differ from dedicated ESG or sustainability reports. \\ \hline
\textbf{George \& Saji} \parencite{s_s__george_b30d4b3e} & \textbf{ESGBench}: explainable question-answering framework. & Domain-grounded questions across multiple ESG themes. & Text \& evidence chains & Created a benchmark requiring supporting evidence for answers, improving traceability. & Does not specifically address deceptive or greenwashed evidence. \\ \hline
\textbf{Schimanski et al.} \parencite{tobias_schimanski_4e192ec5} & \textbf{Multi-domain pretraining}: specialised E, S, and G classifiers. & 13.8M text blocks; three specialised 2k datasets. & Text & Provided targeted datasets for training models that explain variation in ESG ratings. & Final classification datasets are small compared with the pretraining corpus. \\ \hline
\textbf{Schimanski et al.} \parencite{tobias_schimanski_241ada20} & \textbf{Target assessment pipeline}: differentiating pledge types. & Corporate net-zero and reduction-target documents. & Text (targets/ pledges) & Created a benchmark for distinguishing vague pledges from verifiable science-based goals. & Evaluation is restricted to climate targets, omitting social and governance commitments. \\ \hline
\textbf{Pavlova et al.} \parencite{mariya_pavlova_259be499} & \textbf{ESG-FTSE Corpus}: news-based relevance labelling. & News articles and press releases indexed by the FTSE. & Text & Established a corpus for outside-in verification against public sentiment and media evidence. & News data is noisy and requires filtering before matching report-level claims. \\ \hline
\textbf{Schimanski et al.} \parencite{tobias_schimanski_3549ff9f} & \textbf{Nature-related discovery}: transitioning beyond carbon. & TNFD-aligned nature and biodiversity disclosures. & Text & Developed benchmark resources for nature-related disclosures and biodiversity reporting gaps. & Biodiversity reporting standards are still evolving, creating inconsistent ground-truth labels. \\ \hline
\end{tabularx}
\end{table}
\end{landscape}


\begin{landscape}
\begin{table}[p]
\centering
\caption{State-of-the-Art in ESG Benchmark Design}
\label{tab:esg-benchmark-design}
\small
\renewcommand{\arraystretch}{1.2}
\begin{tabularx}{\linewidth}{|p{1.8cm}|p{3cm}|p{3cm}|p{1.8cm}|X|X|}
\hline
Authors & Method / Approach & Datasets & Modalities Involved & Key Contributions & Limitations \\ \hline
\textbf{Bingler et al.}  \parencite{julia_bingler_268445e3} & \textbf{Cherry-Picking Analysis}: Identifying selective reporting in climate disclosures. & 1.6M paragraphs from TCFD-aligned reports. & Text & Technically mapped the gap between "climate-relevant" and "climate-material" disclosures, proving that companies selectively report positive news  \parencite{julia_bingler_268445e3}. & Does not provide a mechanism for automated cross-source validation against external physical risk data. \\ \hline
\textbf{Yu et al.}  \parencite{yinan_yu_e282fd98} & \textbf{climateBUG Framework}: Sector-specific transparency gap analysis. & Bank-specific climate and TCFD reports. & Text \& Financial Data & Identified the "transparency gap" in the financial sector, providing a dataset specifically for carbon-intensive lending disclosures  \parencite{yinan_yu_e282fd98}. & The benchmark is highly tailored to the banking ontology and may not apply to non-financial corporations. \\ \hline
\textbf{Kılınç et al.}  \parencite{yavuz_k_l_n__64e32f80} & \textbf{Deception-Oriented Dataset}: Focused on linguistic indicators of greenwashing. & Multi-sector corporate sustainability reports. & Textual Narratives & Established a dataset to test for "strategic ambiguity" and "vagueness," which are often omitted from standard NLP benchmarks  \parencite{yavuz_k_l_n__64e32f80}. & Performance is limited by the subjective nature of what constitutes "ambiguity" in different regulatory jurisdictions. \\ \hline
\textbf{Billert \& Conrad}  \parencite{fabian_billert_a5fbb687} & \textbf{Nano-ESG Corpus}: External event-based extraction. & News articles covering corporate environmental/social events. & Text & Addressed the "internal-only" data gap by creating a corpus that allows models to verify report claims against public events  \parencite{fabian_billert_a5fbb687}. & Challenges arise in temporal alignment—matching a news event date to the specific reporting cycle of a firm. \\ \hline
\textbf{Wang et al.}  \parencite{xiaojia_wang_02351a60} & \textbf{Sentiment-Thematic Disparity}: Predicting greenwashing degree. & Diverse corporate ESG reports. & Text & Technically quantified the gap between high-sentiment "marketing" language and low-density "action" themes  \parencite{xiaojia_wang_02351a60}. & The model predicts a "degree" of greenwashing but cannot pinpoint specific false statements for auditing purposes. \\ \hline
\textbf{Ong et al.}  \parencite{keane_ong_369a0b63} & \textbf{Aspect-Action Generalization}: Cross-category robustness testing. & ESG reports with labeled aspects and actions. & Text & Developed a benchmark to test how well models generalize "action-detection" across different ESG pillars (E vs. S vs. G)  \parencite{keane_ong_369a0b63}. & Requires highly complex, multi-layered annotations which are difficult to scale for large datasets. \\ \hline
\end{tabularx}
\end{table}
\end{landscape}
     
\chapter{Methodology}
  \label{implementation}
  \section{Research Design and Methodological Overview}

\begin{figure}[ht]
\begin{center}
  \includegraphics*[width=\textwidth]{03_01_overview}
\end{center}
\caption{High-level methodological flow adapted from the existing thesis workflow Mermaid diagrams in the repository.}
\alt{High-level workflow diagram showing the thesis methodology from source reports through OCR, extraction, validation, and thesis-ready evidence outputs.}
\label{fig:methodological_overview}
\end{figure}

This study adopts an executable mixed-method research design for automated ESG disclosure analysis over Indonesian sustainability and annual reports. The methodology is computational because the core workflow converts report PDFs into machine-readable evidence, extracts structured ESG records, classifies aspect-sentiment-tone signals, aligns outputs to an ontology, and evaluates stability across models and prompts. At the same time, it remains interpretive because the workflow preserves page-level provenance, supports manual review through Streamlit audit pages, and treats annotation, disagreement analysis, and failure inspection as part of the research method rather than only a post-processing step.

The central methodological strategy is a staged pipeline implemented in this repository. The pipeline begins with PDF ingestion and OCR, continues with page-aware LLM extraction, adds validation and representation layers through ClimateBERT-style comparison and ABSA-style processing, then produces thesis-ready evidence artifacts such as CSV summaries, JSONL benchmark files, visualizations, and audit tables. In this design, the main inputs are sustainability and annual report PDFs together with associated page markdown, OCR metadata, prompts, and ontology resources. The main outputs are structured ESG records in results/esg\_records.json, stage-specific benchmark outputs such as results/t1\_results.jsonl and results/t2\_results.jsonl, ontology coverage tables, stability summaries, and dashboard exports under results/thesis\_workflow\_dashboard/ and results/revision\_analysis/.

The design intentionally combines more than one methodological logic. First, it uses heuristic and weakly supervised components, including lexicon rules and ontology mappings, because a fully labeled Indonesian ESG corpus is not yet available. Second, it uses multiple LLMs and prompt templates as a comparative extraction strategy, since output quality depends not only on model family but also on prompt format and schema constraints. Third, it uses ClimateBERT-style labeling as an external semantic comparison layer rather than as a direct replacement for disclosure-tone analysis. This separation is important because climate-topic detection and disclosure maturity are related but not identical constructs.

This methodological direction is suitable for the research problem because ESG reports are long, bilingual or code-switched, structurally inconsistent, and often mix narrative claims with metrics, tables, governance boilerplate, and forward-looking commitments. A simple single-model classifier would hide those differences and make source auditing difficult. The proposed workflow instead prioritizes traceability, modularity, and cross-checking. Each page, run, record, and downstream artifact can be traced back to document folders in data/thesis\_dataset/, page markdown files, and logged run metadata. This directly supports the thesis objective of building a tone-aware, ontology-aligned ESG analysis framework for Indonesian reports.

Traditional or simpler alternatives were considered implicitly in the implementation. A pure keyword system would be more deterministic but too brittle for bilingual reporting language. A single LLM prompt would be easier to run but would hide prompt sensitivity and schema drift. A ClimateBERT-only approach would provide climate-topic labels but would not distinguish commitment, action, and outcome tones. The implemented methodology therefore uses comparative prompts, multiple model families, rule-based baselines, classical machine learning, and a lightweight hybrid model as complementary components rather than mutually exclusive replacements.

At the level of research questions, the methodology addresses four linked goals: transforming PDFs into structured ESG evidence, defining aspect-pillar-sentiment-tone schema, comparing tone with ClimateBERT-style labels, and diagnosing instability or failure modes across prompts and models. Reproducibility and provenance are treated as thesis-wide support conditions for those four goals rather than as standalone research questions. Chapter 3 therefore focuses not on a single classifier in isolation, but on a complete methodological system that turns raw disclosures into auditable research evidence.

\subsection{System Architecture}

The repository documentation defines the workflow as an end-to-end research-data pipeline. Its operational architecture can be summarized as follows.

\begin{figure}[ht]
\begin{center}
  \includegraphics*[width=\textwidth]{03_01_01_system_architecture}
\end{center}
\caption{Executable system architecture of the repository, rewritten from the repository's Mermaid workflow documentation into a LaTeX-safe diagram.
  \copyrightstring\ \href{https://creativecommons.org/licenses/by/4.0/}{CC BY 4.0}.}
\alt{System architecture diagram showing the OCR layer, extraction layer, benchmarking layer, and analysis layer connected through repository artifacts and Streamlit pages.}
\label{fig:system_architecture}
\end{figure}

The OCR layer creates page-level markdown and OCR JSON under data/thesis\_dataset/<document>/. The extraction layer uses prompts stored in prompt/ and writes structured outputs and raw responses into results/esg\_records.json and results/background\_llm\_jobs/. The ground-truth and benchmark layer writes resumable JSONL outputs to results/t1\_results.jsonl and results/t2\_results.jsonl. The analysis layer then produces ontology coverage, failure modes, prompt stability, model stability, and greenwashing-oriented summaries in results/revision\_analysis/ and results/thesis\_workflow\_dashboard/.

This architecture is also reflected in the Streamlit surface. pages/Bulk\_OCR.py manages ingestion, pages/llm\_processing.py runs the combined T1-T2-T3 pipeline, pages/ground\_truth.py manages resumable benchmark generation, pages/1\_4\_ClimateBERT\_Record\_Batch.py prepares one-to-one ClimateBERT comparison inputs, pages/1\_6\_Ontology\_Path\_Viewer.py inspects ontology assignments, and pages/1\_0\_Revision\_Analytics.py consolidates diagnostics and stability evidence. The methodology chapter therefore describes a system that is both conceptual and executable.

\subsection{Design Principles}

Five design principles guide the methodology.

First, provenance is preserved at page and record level. Extracted outputs are always linked back to OCR-expanded documents and page markdown, and dedicated verifier pages exist to map structured statements back to their source pages.

Second, the system favors modular comparison over single-model dependence. Prompt families, model families, rule-based logic, classical ML, and hybrid transformer-based components are all retained so disagreement can be inspected rather than hidden.

Third, the pipeline is designed for bilingual robustness. Indonesian and English text may coexist within the same report or even the same segment, so the method uses multilingual preprocessing, bilingual prompts, and ontology resources that reduce cross-language drift.

Fourth, the method prioritizes interpretability together with automation. Lexical triggers, ontology paths, record-level reasoning fields, and audit tables are kept because the thesis is not only measuring output counts but also evaluating whether the outputs are credible and explainable.

Fifth, the workflow is reproducibility-oriented. Background jobs store status files and event logs, prompts are versioned as standalone markdown files, and dashboard outputs are exported into persistent result folders. This principle matters because the thesis relies on repeated reruns, prompt comparisons, and artifact regeneration rather than a single static experiment.

These principles lead directly to later technical decisions: page batching instead of isolated page inference, JSONL resumability for benchmark runs, ontology alignment after extraction, and separate treatment of topic labels, tone labels, and sentiment polarity.

\section{Data Sources}

\begin{figure}[ht]
\begin{center}
  \includegraphics*[width=\textwidth]{03_02_data_sources}
\end{center}
\caption{Section-level diagram for the data-source stack used by the workflow.}
\alt{Data-source diagram showing raw PDF reports, OCR-expanded folders, metadata tables, ontology resources, and downstream analytical artifacts used by the workflow.}
\label{fig:data_source_stack}
\end{figure}

The primary data source is a corpus of sustainability and annual report PDFs stored in data/thesis\_pdf/. According to the repository-wide inspection documented in code\_documentation.md, this directory contained 193 PDF files at the time of inspection. These files form the raw document-level source for the thesis workflow. They are complemented by a machine-expanded corpus in data/thesis\_dataset/, which contained 189 OCR-processed document folders at inspection time. Each processed folder typically includes ocr\_result.json, a pages/ directory with page-level markdown such as page\_0001.md, and an images/ directory with extracted image crops.

This dataset is appropriate for the research problem because the thesis studies how sustainability disclosures can be transformed into structured ESG evidence rather than how short isolated sentences can be classified in a clean benchmark setting. The corpus captures the real reporting environment: long reports, mixed Indonesian and English text, varied layout structure, numeric tables, narrative commitments, governance descriptions, and different sectoral vocabularies. That diversity is necessary for evaluating whether an ESG ABSA pipeline remains useful under realistic disclosure conditions.

The dataset also supports the research questions directly. For RQ1, it provides raw PDFs and OCR-expanded page units. For RQ2 and RQ3, it provides disclosure text that can be converted into aspect, pillar, sentiment, tone, and ClimateBERT-style comparison labels. For RQ4, it provides enough variability to observe schema drift, missing labels, ontology gaps, and model/prompt instability.

The repository includes additional structured support data. data/idx\_data.csv and data/stock\_info/ provide company and sector context. data/ESG Score.xlsx functions as an external benchmark or lookup source rather than as a native pipeline artifact. The main ontology and analysis-side derived resources are stored under results/revision\_analysis/, including ontology.json, ontology coverage tables, prompt stability summaries, failure-mode tables, and pilot annotation files. These resources are not raw data in the same sense as the PDFs, but they are methodological inputs because later stages depend on them for mapping, validation, and evaluation.

\subsection{Technical Characteristics and Sampling}

The raw source data is document-based, but the effective unit of downstream computation changes by stage. OCR produces document folders and page markdown. LLM extraction operates on page batches or page-level contextual chunks. Ground-truth and benchmark stages operate on text units derived from extracted records. Ontology mapping operates on extracted aspects and record-level texts. This multi-level structure is necessary because no single unit is sufficient for all tasks.

The current workflow evidence snapshot  reports 23 completed OCR documents covering about 5,512 pages, 332 extracted tone records, 2,074 T2 rows, 70 pilot labels, and 1,220 result artifacts. These numbers should be interpreted as the active experimental evidence layer used by the thesis dashboard, not as the total raw corpus size. In other words, the repository contains a larger raw and OCR-expanded inventory, while the dashboard snapshot represents the currently processed and thesis-integrated subset.

Selection is therefore partly corpus-driven and partly workflow-driven. Reports are collected because they are relevant to Indonesian ESG disclosure analysis, but the final processed subset is also shaped by practical constraints such as OCR completion, page-batch processing, background job success, and the availability of downstream annotation and validation capacity. This is consistent with the repository design, which stores both large-scale source inventory and smaller thesis-ready evidence layers.

\subsection{Inclusion, Exclusion, and Evidence-Layer Boundaries}

The study uses different inclusion and exclusion logic at different stages, and these boundaries need to be stated explicitly because later evaluation claims depend on them. At document level, the broad inclusion target is Indonesian sustainability or annual reporting material available as machine-processable PDF files in the thesis corpus. Documents are excluded from the active experimental subset when OCR expansion is incomplete, page structure is too unstable for downstream use, quantitative table structure cannot be preserved reliably enough for downstream review, or later extraction and audit stages have not yet produced reusable artifacts.

At page and text-unit level, inclusion is based on usability for provenance-preserving extraction rather than on an assumption that every page contains ESG evidence. Pages may still enter OCR storage even if they later prove non-informative. Table-bearing pages are not treated as a minor formatting variation of narrative pages; they form a special evidence type because they may contain the only explicit quantitative KPI values for a disclosure. For that reason, the workflow distinguishes between narrative-ready pages, table-ready pages, and pages that remain OCR-visible but are too structurally unstable for thesis-facing extraction. In contrast, extracted-record evaluation excludes malformed outputs, explicit missing-tone failures for some downstream denominators, and records that cannot be aligned reliably enough for the specific comparison being reported. This means that the thesis uses several evidence layers rather than a single monolithic dataset: the raw PDF inventory, the OCR-expanded inventory, the active thesis-facing subset, and the smaller reviewed or comparison-ready subset used in specific tables.

The purpose of these boundaries is methodological transparency rather than aggressive filtering. The thesis does not claim that the current active subset is a statistically representative sample of all sustainability reports. Instead, it claims that this subset is sufficient to evaluate whether the implemented workflow can operate on long, noisy, bilingual reports and reveal its main strengths and weaknesses.

\subsection{Ethics, Bias, and Data Limitations}

The corpus is based on corporate reports rather than private personal data, so the main ethical concerns are not individual privacy but rather licensing boundaries, faithful source tracing, and fair interpretation of corporate disclosures. The methodology mitigates evidential misuse by retaining provenance links from structured outputs back to report pages. This reduces the risk that generated records are discussed without access to their source context.

Several biases remain. First, the corpus is biased toward firms that produce machine-accessible reports and toward sectors with stronger reporting practices. Second, environmental and governance language appears more strongly represented than social disclosures in the current evidence snapshot, which may affect downstream tone and aspect distributions. Third, bilingual and code-switched language can produce uneven extraction quality because the same disclosure function may be expressed differently across Indonesian and English segments. Fourth, prompt-dependent extraction can introduce representation bias if one template systematically omits certain fields.

Data quality limitations are explicitly acknowledged in the repository. OCR noise, inconsistent formatting, page duplication, table fragmentation, and missing or unstable JSON fields can all affect downstream outputs. The thesis dashboard also notes a current need for stronger OCR baselines and formal manual validation. These limitations do not invalidate the workflow, but they do affect generalizability and motivate the use of diagnostics, review queues, and pilot annotation rather than overclaiming final benchmark performance.

\section{Data Collection and Preprocessing Pipeline}

\begin{figure}[ht]
\begin{center}
  \includegraphics*[width=\textwidth]{03_02_02}
\end{center}
\caption{Mermaid-derived preprocessing sequence from PDF ingestion to auditable ESG-ready text units.}
\alt{Preprocessing sequence diagram showing PDF ingestion, OCR expansion, page-level markdown generation, storage, and conversion into extraction-ready text units.}
\label{fig:preprocessing_sequence}
\end{figure}

The preprocessing pipeline converts raw PDFs into analyzable page units and then into extraction-ready text batches. This stage is methodologically important because every later classifier or ontology mapper depends on the quality and structure of these intermediate artifacts.

\subsection{OCR Expansion and First-Stage Extraction}

The first transformation is document OCR and structural extraction. The repository documents this stage through pages/Bulk\_OCR.py, pages/1\_2\_OCR\_Quality\_Workbench.py, and the OCR-expanded folders in data/thesis\_dataset/. Each processed report yields an ocr\_result.json file containing page arrays and fields such as markdown text, images, tables, hyperlinks, dimensions, and confidence-related metadata. In parallel, the system writes page-level markdown files under pages/, which are the most commonly reused unit in downstream page audits and page-batch inference.

The unit of analysis at this point is the page. This choice is appropriate because later provenance checks require the workflow to map structured ESG statements back to specific page files, and page-level markdown provides a stable and auditable representation of the original report. At the same time, the system does not restrict itself to single-page inference; page-level files can be grouped into page batches when additional context is needed.

\subsection{Table Extraction and Quantitative Evidence Preservation}

Quantitative ESG evidence cannot be assumed to survive if the workflow treats all pages as plain narrative text. Sustainability reports frequently place emissions, waste, training, governance, and safety indicators inside dense tables, sometimes split across several pages. If those structures are flattened too early, the pipeline may preserve keywords while losing the row, column, unit, or year relationship that gives the number its meaning.

For that reason, the revised methodology treats table-bearing pages as a dedicated ingestion branch rather than as an afterthought of OCR. During OCR expansion, detected table regions are preserved together with their page identifiers, geometric context, and surrounding markdown. When a table spans multiple pages, header continuity, repeated stub columns, and metric-unit consistency are used as reconstruction cues. The intended output of this branch is not merely table text, but row-level quantitative evidence objects that retain document id, page id, table id, row content, normalized metric values when possible, and a provenance pointer back to the source pages.

This design changes the evidence boundary stated earlier in the chapter. Narrative extraction and table extraction are treated as parallel but converging paths. Narrative pages remain suitable for page-batch LLM extraction, while table-preserving outputs support KPI-sensitive review and later fusion into the record layer. In other words, the workflow aims to avoid a false choice between interpretive narrative evidence and quantitative tabular evidence. Both are needed for Indonesian ESG analysis because companies often mix commitments in prose with performance metrics in tables.

The current repository already stores table-aware OCR fields, but the revised thesis positions full table reconstruction as an explicit methodological requirement rather than an implicit future convenience. This is important for validity: if quantitative tables are lost, downstream tone or greenwashing interpretations may overstate promises and undercount realized outcomes. The table-preservation branch therefore functions as a safeguard against narrative bias in the evidence layer.

\subsection{Standardization, Storage, and Tooling}

The preprocessing outputs are standardized around a reproducible directory structure. Raw source PDFs remain in data/thesis\_pdf/. OCR-expanded artifacts are stored in data/thesis\_dataset/<document>/ocr\_result.json, data/thesis\_dataset/<document>/pages/, and data/thesis\_dataset/<document>/images/. This structure allows downstream workers to resolve a document folder, inspect its pages, and trace extraction failures to specific files.

The methodological logic is also standardized in code. pages/llm\_processing.py defines repository constants for the prompt directory, OCR output directory, results directory, and background-job directory. Background jobs are launched through code/llm\_background\_worker.py, while run state is preserved through status.json, control.json, and JSONL-like event logs under results/background\_llm\_jobs/. This storage logic is not incidental implementation detail; it is part of the reproducibility design because it ensures that extraction can be resumed, audited, and compared across runs.

\subsection{Conversion into Structured ESG Records}

After OCR, the next conversion is from page text into structured ESG records. In pages/llm\_processing.py, the T3 extraction stage sends contextualized text to LLM backends and appends results to results/esg\_records.json. Each run object can store metadata such as timestamp, model, target pages, prompt, success status, parsed records, raw output, and error information. When successful, record-level fields may include text, aspect, labels, esg, tone, sentiment, sentiment\_score, and reasoning. In the revised methodological framing, this conversion is extended conceptually to cover both narrative-derived records and table-derived evidence records, even when those two branches require different preprocessing before they can be merged.

This conversion stage is central to the methodology because it creates the canonical evidence layer used by the rest of the system. Instead of treating extracted text as an unstructured blob, the workflow transforms it into schema-bearing records that can later be benchmarked, filtered, aligned to ontology paths, or compared against ClimateBERT outputs. The important revision here is that a usable ESG record may originate either from narrative prose or from a reconstructed table row. Both origins must remain visible in the provenance fields so that later analysis can distinguish between claim-like language and metric-bearing evidence.

\subsection{Alignment and Synchronization}

The workflow combines several data layers that must be aligned after extraction. First, extracted records must remain linked to their source document and source pages. Second, T1 and T2 benchmark outputs must be aligned back to the record texts or labels they were derived from. Third, ontology mapping must connect extracted aspects to canonical ontology paths. Fourth, external comparison labels such as ClimateBERT-style outputs must preserve stable record identifiers for later agreement analysis.

This alignment logic is visible in multiple parts of the repository. pages/2\_2\_LLM\_Statement\_Page\_Verifier.py maps extracted statements back to OCR page markdown. pages/1\_4\_ClimateBERT\_Record\_Batch.py explicitly preserves record\_id for one-to-one ClimateBERT comparison runs. code/data\_alignment.py implements matching utilities for aligning reference labels, ABSA outputs, and benchmark outputs by text or related identifiers. The alignment unit therefore varies by task, but the underlying principle is stable: every derived label should remain recoverable to a concrete upstream artifact.

\subsection{Preprocessing Quality Checks}

The repository includes both manual and automated quality checks. Automated checks appear in OCR processing summaries, page processing audits, parse-success tables, and missing-field diagnostics. Manual checks are supported through Streamlit inspection pages such as the OCR quality workbench, page-level processing audit, statement-to-page verifier, and ground-truth review visualizers.

The dashboard report indicates that OCR ingestion is operationally complete for the active subset, but it also notes the need for stronger semantic baselines such as page-count parity checks and small manually reviewed OCR samples. This is methodologically important: preprocessing quality is treated as a measured risk rather than an invisible assumption. The final preprocessed dataset is therefore not just the collection of OCR files, but the subset that passes through enough audit steps to be usable for extraction and benchmarking.

\section{Feature Extraction and Representation Learning}

\begin{figure}[ht]
\begin{center}
  \includegraphics*[width=\textwidth]{03_04}
\end{center}
\caption{Section-level diagram summarizing how explicit, sparse, contextual, and ontology-aware features are combined.}
\alt{Feature-extraction diagram showing lexical rules, TF-IDF features, contextual embeddings, and ontology-aware representations feeding into the ESG analysis workflow.}
\label{fig:feature_strategy}
\end{figure}

Feature extraction in this study is multi-layered. Some features are explicit and interpretable, such as lexical triggers, aspect keywords, and ontology paths. Others are learned or semi-learned, such as TF-IDF vectors, transformer embeddings, and hybrid fused representations. The repository does not rely on one single feature philosophy; instead, it combines surface features and contextual features because ESG report language contains both formulaic disclosure patterns and semantically nuanced claims.

\subsection{Overall Feature Strategy}

The extracted feature groups can be divided into four categories. The first group contains lexical and rule-based indicators used for aspect, polarity, and tone heuristics. The second group contains classical sparse text representations built from word and character n-gram TF-IDF. The third group contains contextual multilingual sentence embeddings and section-aware representations. The fourth group contains metadata-like supporting features such as ontology paths, source section labels, page linkage, and comparison labels from ClimateBERT-style inference.

These feature groups are connected directly to the research objective. Lexical features improve interpretability and provide a transparent baseline. Sparse features support classical machine learning comparisons. Contextual embeddings capture semantic variation across Indonesian, English, and mixed disclosures. Ontology and metadata features connect extracted outputs to ESG-specific interpretive structure and later evaluation tasks.

\subsection{Rule-Based Lexical Features}

The rule-based component is implemented in code/rule\_based.py with supporting lexicons in code/lexicons.py. The method matches disclosure text against predefined aspect lexicons and tone/polarity trigger lists. collect\_aspects() assigns one or more aspect labels by matching regular-expression patterns. polarity\_basic() uses positive and negative word lists to assign a simple polarity label. tone\_basic() uses ordered lexical cues to assign disclosure tone, prioritizing Outcome over Action over Commitment, and returning Unknown when no strong cue is found.

This feature group is highly interpretable because every decision can be tied to a lexical trigger. The helper explain\_rule\_based\_sentence() returns matched aspect, sentiment, and tone triggers for a sentence, which is useful for explainability outputs in later chapters. The limitation is that rule-based matching can be brittle, especially under bilingual phrasing, code-switching, implicit claims, or complex sentence structure.

\subsection{Classical Sparse Text Features}

The classical ML component is implemented in code/classical\_ml.py. It builds a Featureizer that combines word n-gram TF-IDF with character n-gram TF-IDF. Word n-grams capture lexical and short-phrase disclosure patterns, while character n-grams add robustness to spelling variation, morphology, and OCR noise. After vectorization, the method trains one-vs-rest logistic regression for multi-label aspects and logistic regression or dummy fallbacks for sentiment and tone.

Formally, if a sentence \(x\) is transformed into a sparse vector \(v(x)\),
the linear classifier computes scores of the form

\[
z_k = w_k^\top v(x) + b_k
\]

for class \(k\), where \(w_k\) is the learned coefficient vector and \(b_k\)
is the intercept. 

The predicted class is derived from the highest or thresholded score depending on the task. This representation is appropriate because it provides a strong, reproducible baseline and makes coefficient-based explanations possible. The code also includes local explanation logic that multiplies active feature values by model coefficients to show which terms contributed most strongly to a prediction.

\subsection{Contextual and Hybrid Representations}

The contextual representation layer is implemented in code/hybrid\_model.py. This module uses a small multilingual transformer encoder, defaulting to distilbert-base-multilingual-cased when available, and falling back to deterministic embeddings if the transformer cannot be loaded. Sentence embeddings are pooled from transformer outputs, then passed into a HierarchicalEncoder that incorporates section-type embeddings and cross-section attention. This allows the method to capture not only the local sentence meaning but also document-structure context.

The hybrid model then fuses four information sources: sentence vectors,
contextual section vectors, ontology vectors, and a document vector. In the
\texttt{MTLHybrid} class, the fused representation is formed by concatenation
followed by a nonlinear projection:

\[
h_i = f\left([s_i ; c_i ; o_i ; d]\right)
\]

where \(s_i\) is the sentence vector, \(c_i\) is the contextual section vector,
\(o_i\) is the ontology vector, \(d\) is the document-level vector, and
\(f(\cdot)\) is the learned fusion network. Separate output heads then predict
sentiment and tone. This representation is appropriate because ESG disclosure
meaning is often shaped by both local phrasing and broader section context,
for example when a metric appears inside a target-setting section or a
governance compliance section.

\subsection{Ontology-Based Feature Enhancement}

Ontology alignment is used as a feature enhancement and interpretive layer rather than only as a final visualization step. In the rule-based pipeline, canonical ontology paths are attached using CANON\_PATHS. In the hybrid pipeline, ontology nodes are embedded and projected into the fused representation. The workflow documentation also describes ontology resources as including pillar, aspect, sub-aspect, synonym, and regulatory-anchor information.

This enhancement improves the representation in two ways. First, it regularizes semantically similar disclosures under common ESG concepts. Second, it makes outputs regulator-readable and more interpretable by mapping extracted aspects to structured paths rather than leaving them as free-text labels only. The main limitation is ontology scope: if an Indonesian-specific disclosure topic is absent from the current ontology resources, the mapping may still be incomplete or overly generic.

\section{Proposed Framework}

\begin{figure}[ht]
\begin{center}
  \includegraphics*[width=\textwidth]{03_05}
\end{center}
\caption{Conceptual split between record generation and validation-comparison within the proposed framework.}
\alt{Framework diagram separating LLM-centered ESG record generation from the downstream validation, benchmarking, and comparison layers.}
\label{fig:framework_split}
\end{figure}

The proposed framework is not a single standalone classifier but a composite methodological system centered on a structured ESG evidence store. The extracted features described above are used in two main ways: first, to generate structured records from report text; second, to evaluate and refine those records through comparison, benchmarking, and ontology-aware analysis.

\subsection{Framework 1: LLM-Centered ESG Record Generation}

The first framework is the page-aware LLM extraction workflow implemented primarily in pages/llm\_processing.py and code/llm\_background\_worker.py. Its purpose is to convert OCR-expanded report text into structured ESG records with fields such as text span, aspect, pillar, tone, sentiment, and reasoning. The framework accepts document or page-batch inputs, combines them with prompt templates, sends requests to configured model backends, and writes normalized outputs into results/esg\_records.json.

The implemented prompt inventory shows multiple prompt families, including zero-shot, few-shot, chain-of-thought, and tone-specific variants in both Indonesian and English, plus data-oriented prompt files such as data.md and data\_v1.md. This comparative prompt design allows the framework to test whether output quality depends on schema framing, language, or reasoning style.

At a high level, the framework processes a batch
\(B = \{p_1, \dots, p_n\}\) of page texts under prompt template \(q\)
and model \(m\) to produce a structured output set \(R\):

\[
R = F(B, q, m)
\]

where \(F\) denotes the end-to-end extraction function. In practice,
\(R\) is accepted only if it can be parsed into the expected record
schema. Otherwise, the run is preserved with error metadata for later
audit rather than silently discarded.

\subsection{Framework 2: Benchmarking, Comparison, and Evidence Scoring}

The second framework is the validation and comparison layer built around T1 and T2 outputs, ClimateBERT-style comparison, and ontology-aligned diagnostics. pages/ground\_truth.py loads extracted texts from results/esg\_records.json, derives benchmark text units, and writes resumable outputs to results/t1\_results.jsonl and results/t2\_results.jsonl. T1 focuses on ClimateBERT or local classification-style outputs, while T2 captures rule-based or hybrid ABSA-oriented labels and related processing.

The design uses JSONL because each processed item can be appended independently and resume logic can skip already complete records. load\_processed\_t1() and load\_processed\_t2() use stable keys such as label-model pairs or labels alone to avoid rerunning completed items. This matters methodologically because benchmark generation may be interrupted by model availability, long-running jobs, or partial failures, and the study needs a reproducible way to continue from the last valid state.

Evidence from multiple sources is then combined downstream. The workflow compares tone labels against ClimateBERT-style labels, maps extracted aspects to ontology paths, measures missing-tone and schema-drift rates, and preserves disagreement cases for manual review. The scoring logic is therefore not a single scalar objective but a collection of evidence-quality indicators: parse success, field completion, label agreement, ontology coverage, and failure-mode distribution.

\subsection{Inference and Output Mapping}

Inference in the overall system means more than generating a label. It includes mapping model outputs back into actionable thesis artifacts. For T3 extraction, raw outputs are parsed into structured records. For T1 and T2 benchmarking, JSONL outputs are mapped back to labels and record texts. For ontology analysis, extracted aspects are mapped to canonical paths. For provenance checks, statements are mapped back to page markdown. For dashboard reporting, all of these outputs are aggregated into CSV summaries, graphs, and narrative evidence tables.

This mapping step is essential because the thesis aims to produce usable disclosure evidence rather than isolated classification scores. A predicted tone label has limited value if it cannot be tied back to a record, source page, ontology path, and comparison context. The implemented framework therefore treats output mapping as part of the methodological core.

\section{Ground Truth, Weak Labels, and Evaluation Reference}

\begin{figure}[ht]
\begin{center}
  \includegraphics*[width=\textwidth]{03_06}
\end{center}
\caption{Layered evaluation-reference construction used when a full expert gold corpus is not yet available.}
\alt{Reference-construction diagram showing extracted records, weak labels, pilot human annotations, and ClimateBERT-style comparison signals combined into a layered evaluation framework.}
\label{fig:reference_construction}
\end{figure}

True expert-labeled reference data does not yet exist for the full corpus. This is one of the central methodological constraints of the study, and the repository is designed around that constraint rather than ignoring it. The workflow therefore uses a layered reference strategy: weak labels from extraction outputs, benchmark outputs in JSONL form, pilot human annotation files, and comparison labels from ClimateBERT-style runs.

\subsection{Why a Special Reference Construction Process Is Needed}

Existing outputs are not sufficient to act as final gold labels. The extracted record layer is generated by prompts and models that can drift in schema, omit fields, or disagree across templates. ClimateBERT-style labels are useful as an external comparison signal but do not fully capture disclosure tone. Human annotation is available only in pilot form. As a result, the study requires a special reference-construction process that balances automation with manual oversight.

The dashboard snapshot confirms this limitation. It reports 70 pilot labels and explicitly notes the need for stronger human-coded baselines, disagreement examples, and broader evaluation coverage. The methodology therefore treats reference construction as iterative: weak labels support early experimentation, pilot labels support focused validation, and future expansion can move toward a stronger stratified gold set.

\subsection{Reference Generation Procedure}

Reference generation begins with extracted records in results/esg\_records.json. pages/ground\_truth.py converts top-level text fields, nested record texts, and parseable raw outputs into benchmark text units. These units are then processed by T1 and T2 pipelines, with outputs appended line-by-line into results/t1\_results.jsonl and results/t2\_results.jsonl. The use of JSONL makes the process resumable and auditable.

For T1, the system can call a ClimateBERT client when available or use local model alternatives discovered from the model cache. For T2, rule-based and hybrid logic provide tone, sentiment, and ontology-related outputs. In addition, pages/1\_4\_ClimateBERT\_Record\_Batch.py prepares one-to-one record batches for external ClimateBERT validation, explicitly preserving record\_id so imported outputs can be aligned back to the original records.

Human review enters through the pilot annotation files and review-oriented Streamlit pages. results/revision\_analysis/pilot\_ground\_truth\_seed.csv provides an annotation scaffold, while pilot\_ground\_truth\_annotations.csv stores reviewed labels. These artifacts do not yet form a complete gold corpus, but they are sufficient to support targeted inspection, early agreement analysis, and the identification of difficult cases.

\subsection{Validation, Constraints, and Current Limitations}

Generated references are validated through a combination of structural and semantic checks. Structurally, a record is considered complete only if expected labels are present and no explicit error state is stored. Semantically, agreement pages and review queues inspect whether tone, sentiment, and ClimateBERT-style labels are coherent or contradictory. Provenance pages further constrain the outputs by checking whether extracted statements can be found in the OCR source text.

The current reference layer nevertheless has important limitations. Weak labels remain model-dependent. Prompt sensitivity affects which records are extracted in the first place. OCR quality is operationally tracked but not yet fully benchmarked with formal CER or WER across the active subset. Ontology coverage, while strong in the current dashboard snapshot, still depends on the scope of the existing ontology resources. For these reasons, the study does not claim a finalized gold-standard benchmark. Instead, it claims an auditable and extensible reference-construction process suitable for iterative thesis research.

\section{Threats to Validity and Reproducibility Boundaries}

Internal validity is limited by the evolving nature of the research workspace. Prompt refinement, model comparison, ontology expansion, and review-oriented diagnostics were developed iteratively in the same repository. This improves engineering transparency, but it also means the thesis should not be read as a fully frozen benchmark in which development and evaluation were perfectly isolated from the beginning.

Construct validity is limited by the fact that disclosure tone, sentiment polarity, climate-topic labels, and greenwashing-oriented ratios are related but non-identical constructs. The thesis mitigates this by keeping those layers separate analytically, but it cannot eliminate the underlying conceptual overlap entirely. ClimateBERT-style commitment outputs are therefore treated as an external semantic comparison signal rather than as definitive reference truth for tone.

External validity is limited by domain concentration. The active evidence layer is drawn from Indonesian sustainability and annual reports, often with bilingual or mixed-language phrasing and sector-specific reporting conventions. The findings therefore support claims about this setting more strongly than claims about other regulatory, linguistic, or sectoral environments.

Reproducibility is strong at the level of stored artifacts, prompts, logs, and resumable outputs, but weaker at the level of exact semantic reruns for third-party LLMs. Model availability, provider-side updates, and stochastic generation behavior can alter future outputs even when the repository preserves prompt files and artifact paths. For that reason, the thesis claims reproducible workflow structure and reproducible evidence storage more confidently than perfectly identical future outputs.

\section{Summary}

\begin{figure}[ht]
\begin{center}
  \includegraphics*[width=\textwidth]{03_07_summary}
\end{center}
\caption{Summary diagram of the chapter's end-to-end methodological contribution.}
\alt{Summary diagram showing the end-to-end methodological contribution from report ingestion through structured ESG evidence generation, validation, and thesis reporting.}
\label{fig:methodology_summary}
\end{figure}

This chapter has described an executable methodology for Indonesian ESG disclosure analysis. The workflow begins with PDF reports, expands them into OCR-based page artifacts, transforms those artifacts into structured ESG records through comparative LLM prompting, enriches the records through rule-based, classical ML, and hybrid contextual representations, aligns them to ontology paths and ClimateBERT-style comparison labels, and preserves the full process in reproducible result stores and dashboard artifacts.

Methodologically, the contribution of this design is not only that it automates extraction, but that it does so in a way that remains auditable, modular, bilingual-aware, and compatible with partial reference data. The next chapter can therefore focus on implementation and results using a clearly defined data flow, model structure, and evaluation reference layer grounded in the code and artifacts of the repository.























































































































































































































     
\chapter{Experiments}
  \label{experiments}
  \section{Experimental Scope and Protocol}

\subsection{Experimental Objective and Tested Methods}

The purpose of the Chapter 4 experiments is to evaluate whether the implemented ESG pipeline can transform sustainability-report PDFs into structured, auditable ESG disclosure records, classify those records into ESG pillar, aspect, sentiment, and disclosure tone, compare tone outputs against ClimateBERT-style labels, and expose failure modes that matter for thesis validity. In practice, the experiments do not test one isolated model. They test a full repository workflow composed of OCR ingestion, LLM-based structured extraction, benchmark labeling, ontology alignment, and diagnostic dashboards.

The main method under evaluation is the repository’s staged ESG analysis framework. The first stage is OCR and page extraction, which converts PDF files into page-level markdown and OCR JSON. The second stage is T3 LLM extraction, which generates structured ESG records stored in results/esg\_records.json. The third stage is the T1 comparison layer, where extracted record texts are compared against ClimateBERT-style classification outputs or their local proxy equivalents. The fourth stage is the T2 ABSA-style layer, which includes rule-based logic, classical machine learning, and a lightweight hybrid contextual model. The final layer aggregates outputs into ontology coverage tables, failure-mode summaries, agreement scores, and dashboard artifacts.

Several comparison approaches are already implemented in the codebase. code/rule\_based.py provides a lexicon-driven baseline with explicit aspect, polarity, and tone triggers. code/classical\_ml.py provides a classical TF-IDF plus logistic-regression baseline. code/hybrid\_model.py provides a contextual multilingual hybrid representation that fuses sentence embeddings, section context, ontology vectors, and a document-level vector. In parallel, the LLM extraction layer compares prompt and model variants rather than relying on a single fixed prompt. This is important because the thesis problem is not only classification performance but also schema stability, field completion, and practical utility of extracted records.

Because the corpus is centered on Indonesian disclosures, multilingual coverage alone is not sufficient as an evaluation logic. A revised comparison condition is therefore required: wherever possible, the multilingual architecture should be read against an Indonesian-oriented baseline rather than being treated as self-justifying for the local language context. In this thesis, that requirement is framed as a necessary baseline extension for the reviewed subset rather than as a claim that the current saved dashboard already contains a complete IndoBERT benchmark.

The implemented prompt inventory includes seven primary prompt templates used in the thesis-facing stability analysis: data.md, tone\_chain\_of\_thought\_english.md, tone\_chain\_of\_thought\_indonesian.md, tone\_few\_shot\_english.md, tone\_few\_shot\_indonesian.md, tone\_zero\_shot\_english.md, and tone\_zero\_shot\_indonesian.md. These operationalize three prompting strategies: zero-shot, few-shot, and chain-of-thought, in both English and Indonesian. At the backend level, the code supports OpenRouter-hosted models, LM Studio or other OpenAI-compatible local endpoints, and Ollama-style local inference. The active results summarized in the thesis dashboard show that the most important model comparison for the current evidence layer is between arcee-ai/trinity-large-preview:free and openai/gpt-oss-120b:free, while additional live reprocess runs include arcee-ai/trinity-large-thinking:free, minimax/minimax-m2.5:free, and openai/gpt-oss-20b:free.

This design follows directly from Chapter 3. The methodology argued that a mixed and modular evaluation strategy is more appropriate than a single-model setup because ESG reports are bilingual, structurally inconsistent, and prone to OCR and schema noise. The Chapter 4 implementation therefore tests both final output quality and engineering reliability.

\subsection{Configuration, Infrastructure, and Reproducibility}

The implementation is centered on a Python and Streamlit research workspace. The main system surface is a multi-page Streamlit application under pages/, supported by reusable logic in code/. Source data is stored under data/, derived artifacts under results/, prompt templates under prompt/, and documentation under documentation/. This structure is part of the experiment design because each stage writes stable intermediate artifacts that can be audited or reused later.

For OCR-expanded data, each processed document in data/thesis\_dataset/ contains ocr\_result.json, page markdown files, and extracted images. For extraction outputs, pages/llm\_processing.py writes structured run objects to results/esg\_records.json and background job state to results/background\_llm\_jobs/. For benchmark runs, pages/ground\_truth.py writes resumable JSONL outputs to results/t1\_results.jsonl and results/t2\_results.jsonl. For thesis reporting, summary tables and figures are exported to results/revision\_analysis/ and results/thesis\_workflow\_dashboard/.

The main feature-extraction and model settings are as follows. The rule-based model uses curated lexical triggers for aspect, polarity, and tone. The classical model uses word and character TF-IDF vectors with logistic regression and one-vs-rest classification for multi-label aspects. The hybrid model uses distilbert-base-multilingual-cased when available, with a lightweight hierarchical encoder and a fused prediction head. In code/hybrid\_model.py, the encoder weights are frozen by default for fast CPU-friendly operation, and the architecture uses a small multi-head attention block for section interaction. This is appropriate for the current thesis environment because the repository is designed to be executable on local research hardware and interactive dashboards, not only on high-end training servers.

The experiment results also reflect the available compute environment. The active pipeline supports cloud LLM calls through OpenRouter and local model execution through LM Studio-compatible backends and downloaded Hugging Face-style models. ClimateBERT comparison in the current saved outputs mostly uses local downloaded models such as distilroberta-base-climate-commitment and related classification checkpoints, while the repository notes that a full one-to-one remote ClimateBERT benchmark is still future work. The same boundary applies to Indonesian-monolingual benchmarking: the revised manuscript treats an IndoBERT-style baseline as a required local-language comparison condition, but does not misrepresent the current saved evidence as if a full fine-tuned benchmark were already complete. This constraint affects how results are interpreted: current ClimateBERT tables are meaningful as proxy validation, but not yet as a final external benchmark, and local-language benchmarking remains an explicit extension target for stronger validity.

Reproducibility is one of the strongest implementation features of the codebase. The thesis dashboard report lists 1,220 result artifacts, 184 LLM background jobs, and persistent saved graph attachments. The repository explicitly states that visualizations can be regenerated by rerunning scripts such as code/visualize\_tone\_climatebert.py, while Streamlit pages such as 6\_1\_Chapter\_4\_Implementation\_Results.py and 6\_6\_Chapter\_4\_Results\_Visualizer.py provide live views over the same stored artifacts. The experiments are therefore reproducible in an engineering sense even where some semantic baselines remain incomplete.

\subsection{Evaluation Boundaries and Comparison Conditions}

Not all Chapter 4 results use the same denominator or the same evidence layer, so each comparison must be interpreted against its own subset. OCR completion refers to the tracked processed-document layer. Tone distributions and agreement tables refer to the active extracted-record layer after additional validity filtering. Pilot annotation observations refer to a smaller review-oriented subset. For that reason, the chapter avoids treating every reported number as if it came from one unified benchmark corpus.

Comparison fairness is strongest when the same extracted text unit is preserved across later stages. This is the case for the main tone-versus-ClimateBERT comparison workflow, where record-level text and record identifiers are preserved for downstream comparison. By contrast, prompt and model comparisons should be interpreted as controlled configuration comparisons only when the same data slice and the same acceptance rules are held constant. Where those conditions are not perfectly fixed, the thesis treats the results as diagnostic evidence rather than as definitive head-to-head benchmark rankings.

\section{Evaluation Metrics}

\subsection{Primary Metrics}

The thesis uses several primary evaluation metrics because no single metric captures the full quality of this pipeline.

For RQ1, the main metric is OCR processing completion, measured as the number of documents successfully processed and the number of pages covered by the OCR-expanded corpus. In the current dashboard snapshot, 23 out of 23 OCR-tracked documents are marked done, covering approximately 5,512 pages. This metric is suitable because the first requirement of the system is operational: the pipeline must convert PDFs into usable intermediate artifacts before any downstream ESG analysis can occur.

For T3 extraction quality, the primary metrics are JSON parse success rate, average records per run, field completion rate, missing-tone rate, and schema-drift rate. JSON parse success rate measures whether a model and prompt combination produces structurally valid outputs. Average records per run measures extraction yield. Field completion rate measures how often the expected schema is actually filled. Missing-tone rate measures how often tone is omitted, which is critical because tone is the main thesis contribution. Schema-drift rate captures malformed or repurposed fields, such as tone values appearing in sentiment slots.

For classification and comparison tasks, percent agreement and Cohen’s kappa are used. Agreement captures raw overlap between the repository’s tone labels and ClimateBERT-style commitment labels, while Cohen’s kappa adjusts for chance agreement. In the saved summary, the comparison tone\_commitment\_vs\_climate\_commitment\_label covers 332 records and yields 83.7\% agreement with Cohen’s kappa of 0.645. Higher values indicate stronger alignment, but kappa is preferred over raw agreement alone because the commitment class is relatively frequent.

For representation and ontology quality, ontology coverage is used. This metric tracks whether extracted aspects can be mapped to ontology paths. In the current RQ4 summary, 52 aspects are tracked and all 52 are mapped to ontology paths. This does not prove perfect semantic correctness, but it measures whether the ontology layer is operational and comprehensive enough for downstream interpretation.

For company-level interpretive risk, the greenwashing index is used as a task-specific ratio between commitment and outcome records at document level. Higher values indicate a stronger imbalance toward promises relative to reported outcomes. The metric is useful because it translates the tone taxonomy into an interpretable governance and disclosure-risk signal, but it remains heuristic and is not yet externally validated.

\subsection{Secondary and Complementary Metrics}

Basic metrics alone are insufficient because this pipeline is not a single-label classification benchmark. Additional metrics are needed to evaluate reliability, interpretability, and usefulness.

First, prompt stability metrics are used to assess how sensitive the extraction layer is to prompt formulation. These include prompt-level parse success, average records, missing-tone rate, schema-drift rate, and field completion. Second, model stability metrics assess whether apparent quality gains are robust across model families. Third, failure-mode counts are used to identify specific structural or linguistic conditions under which the system breaks down, including bilingual or code-switched text, modal or hedged language, passive constructions, regulatory Indonesian terms, table-heavy layouts, and explicit missing-tone cases.

Fourth, denominator audits are used to interpret results correctly. The file results/revision\_analysis/chapter4\_tone\_denominator\_audit.csv distinguishes corpus extraction coverage from tone-table denominators. For example, 5,444 units are used for total extracted-record coverage, but 4,853 units are used for tone distribution and agreement because 591 cases are excluded from those denominators and treated as missing-tone quality issues. This is methodologically important because agreement should not be inflated or distorted by invalid or absent tone outputs.

Fifth, lexical trigger counts serve as a lightweight explainability metric. They show how often action, commitment, outcome, hedge, passive, regulatory, and table-layout triggers co-occur with predicted tones. These metrics do not replace semantic evaluation, but they help explain why certain prompts or models tend to overpredict commitment or underdetect outcome language.

Taken together, these metrics provide a more complete evaluation frame. OCR metrics measure operational readiness. Parse and completion metrics measure extraction usability. Agreement metrics measure construct overlap. Ontology coverage measures semantic mapping readiness. Stability metrics measure robustness. Failure modes measure weakness concentration. Greenwashing ratios and lexical triggers provide task-specific interpretation beyond generic accuracy.

\subsection{Local-Language Baseline Condition}

An additional evaluation condition is required by the thesis domain itself: the workflow should be interpretable relative to a local-language baseline. Indonesian ESG disclosures contain governance formulae, regulatory terminology, and company-specific phrasing that may be underrepresented in general multilingual encoders. For that reason, a fine-tuned Indonesian-monolingual baseline such as IndoBERT is the most appropriate comparison target for the reviewed subset when the task is sentence- or record-level tone discrimination.

The purpose of this baseline is not to replace the multilingual architecture outright. It is to test whether Indonesian-specialized representations improve performance in cases that are currently difficult for the pipeline, especially governance-heavy statements, regulation-sensitive phrases, and mixed promise-action wording. The present chapter therefore interprets the multilingual results cautiously: they are operationally meaningful, but they are not the strongest possible language-specific benchmark until an Indonesian-monolingual comparison has been completed on the same reviewed slice.

\section{Experimental Results}

\subsection{RQ1: ESG ABSA Schema}

RQ1 asks how ESG disclosures can be represented effectively through a record-level schema that integrates aspect, ESG pillar, sentiment, and tone. The Chapter 4 results show that the proposed schema is operational and expressive enough to support that goal in the current thesis-facing subset.

The first result relevant to RQ1 is that the ingestion and extraction pipeline is operational at meaningful corpus scale. The OCR processing summary lists 23 processed documents, all marked as done, covering approximately 5,512 pages. The largest processed reports include Bank Neo Commerce Annual and Sustainable Report Tahun 2024.pdf with 694 pages, vktr\_ar\_sr\_2024.pdf with 548 pages, and alfamart\_sustainability\_2025.pdf with 510 pages. This demonstrates that the schema is not being tested only on short pilot inputs, but on long, structurally varied corporate reports.

The second result is that the extraction layer produces a usable structured evidence set. The thesis dashboard report records 332 tone-bearing ESG records and 2,074 T2 rows. Among the 332 extracted records, the most common tone is commitment with 115 records, and the most common ESG pillar is environmental with 179 records. Governance contributes 121 records, while the social pillar appears only 4 times in the summary used by the thesis planning notes. This indicates that the schema is already expressive enough to separate tone from polarity and from ESG pillar, but it also reveals imbalance in the active sample.

Figure~\ref{fig:tone_distribution} and Figure~\ref{fig:esg_by_tone} make the basic record composition visually transparent before the chapter moves into agreement and failure analysis.

\begin{figure}[H]
\begin{center}
  \includegraphics*[width=\textwidth]{tone_distribution}
\end{center}
\caption{Tone distribution across the active extracted-record layer. This figure supports RQ2 by showing that commitment is the dominant disclosure tone in the current thesis-facing subset.}
\alt{Bar chart showing the distribution of extracted records across tone categories, with commitment as the dominant class.}
\label{fig:tone_distribution}
\end{figure}

\begin{figure}[H]
\begin{center}
  \includegraphics*[width=0.8\textwidth]{esg_by_tone}
\end{center}
\caption{ESG pillar distribution by tone. This figure shows that environmental and governance disclosures dominate the active evidence layer and helps explain why social-pillar findings remain comparatively weak.}
\alt{Stacked bar chart showing ESG pillar counts within each tone category, with environmental and governance records dominating the distribution.}
\label{fig:esg_by_tone}
\end{figure}

Prompt-level extraction behavior is shown in Table~\ref{tab:prompt-level-extraction-performance-across-the-thesis-facing-subset}.

\begin{table}[htbp]
\centering
\caption{Prompt-level extraction performance across the thesis-facing subset.}
\label{tab:prompt-level-extraction-performance-across-the-thesis-facing-subset}
\footnotesize
\setlength{\tabcolsep}{4pt}
\renewcommand{\arraystretch}{1.15}

\begin{tabularx}{\linewidth}{
    >{\raggedright\arraybackslash}p{0.40\linewidth}
    *{6}{>{\centering\arraybackslash}p{0.075\linewidth}}
}
\toprule
\textbf{Prompt}
&
\textbf{Runs}
&
\makecell{\textbf{Parse}\\\textbf{Success}}
&
\makecell{\textbf{Avg.}\\\textbf{Records}}
&
\makecell{\textbf{Missing}\\\textbf{Tone}}
&
\makecell{\textbf{Schema}\\\textbf{Drift}}
&
\makecell{\textbf{Field}\\\textbf{Completion}}
\\
\midrule

\texttt{tone\_chain\_of\_thought\_english.md}
& 16 & 1.000 & 6.25 & 0.0000 & 0.0037 & 0.3744 \\

\texttt{tone\_chain\_of\_thought\_indonesian.md}
& 15 & 1.000 & 4.07 & 0.0028 & 0.0028 & 0.3315 \\

\texttt{tone\_few\_shot\_english.md}
& 15 & 1.000 & 0.00 & 0.0000 & 0.0000 & 0.0000 \\

\texttt{tone\_few\_shot\_indonesian.md}
& 14 & 1.000 & 1.00 & 0.0000 & 0.0000 & 0.0714 \\

\texttt{tone\_zero\_shot\_english.md}
& 14 & 1.000 & 3.93 & 0.0000 & 0.0000 & 0.3571 \\

\texttt{tone\_zero\_shot\_indonesian.md}
& 16 & 1.000 & 2.63 & 0.0000 & 0.0000 & 0.1250 \\

\bottomrule
\end{tabularx}
\end{table}

These results show that parse success alone is not an adequate quality metric. Every prompt listed above reaches 100\% parse success, yet their practical usefulness differs sharply. tone\_chain\_of\_thought\_english.md yields the highest average record count. In other words, a prompt can be syntactically valid and still be unsuitable for the core thesis task.

The denominator audit clarifies the scale of the tone-quality problem. Of 5,444 extracted units in the relevant coverage view, 591 are excluded from tone distribution and agreement because they represent missing-tone cases. This is a meaningful failure rate, not a negligible noise band. It supports the claim that tone extraction is the most fragile part of the schema and justifies treating missing-tone behavior as an explicit Chapter 4 result rather than a minor implementation bug.

Figure~\ref{fig:aspect_by_tone_heatmap} complements the aggregate tone distribution with aspect-level structure.

\begin{figure}[H]
\begin{center}
  \includegraphics*[width=\textwidth]{aspect_by_tone_heatmap}
\end{center}
\caption{Aspect-by-tone heatmap for the active extracted-record layer. This figure supports RQ2 by showing which aspects are concentrated in commitment, action, or outcome language rather than only reporting overall counts.}
\alt{Heatmap showing how extracted ESG aspects are distributed across commitment, action, outcome, and other tone categories.}
\label{fig:aspect_by_tone_heatmap}
\end{figure}

The ontology layer strengthens the schema interpretation further. The current ontology table shows 52 aspects tracked and 52 mapped to ontology paths. This does not prove perfect semantic correctness, but it shows that the record-level schema is not only syntactically structured; it is also semantically linkable to a broader ESG concept layer. In practical terms, the schema supports page-linked extraction, multi-field labeling, and ontology-aligned interpretation in a single evidence structure.

Overall, RQ1 is answered positively within the current thesis scope. ESG disclosures can be represented through a record-level schema that integrates aspect, ESG pillar, sentiment, and tone, and that schema remains usable at realistic report scale. The main qualification is that class balance is uneven and tone completion remains more fragile than aspect or ontology coverage.

\subsection{RQ2: Tone vs. Climate-Specific Models}

RQ2 asks how LLM-generated tone labels compare with specialized outputs from models like ClimateBERT, and what that comparison reveals about the validity of automated ESG assessment. The most important quantitative result for this question is the alignment between disclosure tone and ClimateBERT-style commitment labels. Table~\ref{tab:tone-commitment-versus-climatebert-style-commitment-comparison} summarizes the main agreement result.

\begin{table}[htbp]
\centering
\caption{Tone-commitment versus ClimateBERT-style commitment comparison.}
\label{tab:tone-commitment-versus-climatebert-style-commitment-comparison}
\footnotesize
\setlength{\tabcolsep}{4pt}
\renewcommand{\arraystretch}{1.15}
\begin{tabularx}{\linewidth}{|X|X|X|X|X|X|X|}
\hline
\textbf{Compared labels} & \textbf{Records} & \textbf{Percent agreement} & \textbf{Cohen's kappa} & \textbf{Tone commitment rate} & \textbf{Climate-commitment label rate} & \textbf{Interpretation boundary} \\ \hline
\texttt{tone\_commitment} vs \texttt{climate\_commitment\_label} & 332 & 83.7\% & 0.645 & 34.6\% & 36.4\% & Strong overlap, but not label equivalence \\ \hline
\end{tabularx}
\end{table}

This is a strong but not perfect alignment. The result suggests that the repository’s tone taxonomy captures a signal that overlaps substantially with climate-commitment detection, but the two constructs are not identical. This is theoretically useful because it supports the claim that tone adds a maturity or disclosure-function layer beyond climate-topic recognition.

Figure~\ref{fig:climatebert_label_by_tone} is presented with Table~\ref{tab:tone-commitment-versus-climatebert-style-commitment-comparison} because the crosstab structure is easier to interpret visually than in prose alone.

\begin{figure}[H]
\begin{center}
  \includegraphics*[width=\textwidth]{climatebert_label_by_tone}
\end{center}
\caption{Cross-distribution of LLM tone labels and ClimateBERT-style labels. This figure supports RQ3 by showing where commitment aligns strongly and where action or outcome labels diverge from climate-topic categorization.}
\alt{Stacked chart comparing tone categories against ClimateBERT-style labels to show areas of overlap and divergence.}
\label{fig:climatebert_label_by_tone}
\end{figure}

The saved tone-by-ClimateBERT crosstab provides a more detailed view:

Although the label space is broader than a single climate-commitment binary, the table shows that commitment records cluster differently from action and outcome records. In the thesis planning notes, commitment is most strongly associated with climate-commitment, climate-d, and environmental-claims, while action and outcome disperse more across governance and other label families. This pattern supports the idea that commitment language dominates environmental claim-making even when a record does not yet describe measurable outcome.

The comparison therefore supports a qualified validity claim. Tone labels and ClimateBERT-style labels are strongly related at commitment level, which suggests that the tone taxonomy is not arbitrary. At the same time, the two systems answer different analytical questions. ClimateBERT-style outputs focus on climate-topic or climate-commitment relevance, whereas the tone schema asks whether a disclosure is a promise, an action, or an achieved outcome. The result is therefore best interpreted as partial construct alignment rather than label equivalence.

Overall, RQ2 is answered positively but with a clear boundary. LLM-generated tone labels overlap meaningfully with specialized climate-oriented outputs, especially for commitment language, but they should not be treated as interchangeable with climate-topic classification. The comparison supports the validity of the tone taxonomy as an additional analytical layer rather than as a replacement for specialized climate models.

\subsection{RQ3: Pipeline Diagnostics}

RQ3 asks what specific failure modes characterize the automated extraction of ESG records, including OCR-related data loss, schema problems, and ontology gaps. This is the research question where the Chapter 4 evidence is most explicitly diagnostic rather than comparative.

The failure-mode count table is shown below.

\begin{table}[htbp]
\centering
\caption{Failure-mode count table.}
\label{tab:failure-mode-count-table}
\footnotesize
\setlength{\tabcolsep}{4pt}
\renewcommand{\arraystretch}{1.15}
\begin{tabularx}{\linewidth}{|X|X|X|X|X|}
\hline
\textbf{Failure mode} & \textbf{Count} & \textbf{Share of reviewed failures} & \textbf{Likely cause} & \textbf{Representative example type} \\ \hline
\texttt{missing\_tone} & 61 & 61.0\% & Output omitted the central tone field despite otherwise parseable extraction & Record with populated text/aspect fields but empty tone \\ \hline
\texttt{schema\_drift} & 20 & 20.0\% & Model produced values in the wrong field or changed schema semantics & Tone-like content placed in sentiment or free-text slots \\ \hline
\texttt{hedged\_or\_modal\_language} & 10 & 10.0\% & Commitment/action boundary blurred by future-oriented or hedged phrasing & “will”, “target”, “intend”, “berkomitmen” mixed with current actions \\ \hline
\texttt{regulatory\_or\_indonesian\_domain\_terms} & 3 & 3.0\% & Domain-specific Indonesian ESG terminology not handled consistently & Governance/regulatory phrasing with weak tone cues \\ \hline
\texttt{table\_or\_numeric\_layout} & 3 & 3.0\% & Table-heavy or numeric formatting disrupted semantic extraction & KPI/table rows detached from narrative context \\ \hline
\texttt{passive\_voice} & 3 & 3.0\% & Passive construction weakened action/outcome detection & Statement describes achievement without explicit actor \\ \hline
\texttt{bilingual\_or\_code\_switched} & 1 & 1.0\% & Language switching complicated cue interpretation & Mixed Indonesian-English ESG statement \\ \hline
\texttt{schema\_drift} with \texttt{action} prediction & 1 & 1.0\% & Structural mismatch in otherwise action-like output & Action record with malformed output fields \\ \hline
\texttt{hedged\_or\_modal\_language} with \texttt{action} prediction & 1 & 1.0\% & Action phrasing mixed with promise language & Agreement/plan statement straddling action and commitment \\ \hline
\texttt{passive\_voice} with \texttt{action} prediction & 1 & 1.0\% & Passive form weakened event interpretation & Passive wording treated as action rather than outcome \\ \hline
\end{tabularx}
\end{table}

The dominant weakness is clear: 61 missing-tone cases. Beyond that, schema drift and hedged language are the next most important failure patterns. This means that the pipeline is more vulnerable to discursive ambiguity and formatting irregularity than to simple parser breakdown. The problem is therefore partly linguistic, not only technical.

The revised interpretation of this finding is operational rather than cosmetic. Missing-tone failures should be reduced systematically through stronger schema enforcement, targeted reprompting for empty tone fields, and post-validation against explicit commitment, action, and outcome trigger sets. Records that still fail those checks should remain visible in a review queue rather than being silently discarded. This matters because downstream tone distributions and greenwashing-oriented ratios are sensitive to omission: the system can appear cleaner simply by dropping difficult cases, which would weaken the thesis rather than strengthen it.

Two additional graphs support the failure-analysis argument directly by separating absolute frequency from composition.

\begin{figure}[H]
\begin{center}
  \includegraphics*[width=\textwidth]{failure_mode_pareto}
\end{center}
\caption{Failure-mode Pareto chart. This figure shows which small number of failure categories account for most observed extraction weakness in the current thesis-facing subset.}
\alt{Pareto chart showing failure categories ranked by count with a cumulative-percentage line highlighting the dominant sources of extraction weakness.}
\label{fig:failure_mode_pareto}
\end{figure}

\begin{figure}[H]
\begin{center}
  \includegraphics*[width=\textwidth]{failure_mode_pie}
\end{center}
\caption{Failure-mode composition. This figure shows the relative share of missing tone, schema drift, language ambiguity, layout-related issues, and other recorded failure categories.}
\alt{Pie chart showing the proportional composition of recorded failure categories such as missing tone, schema drift, and language ambiguity.}
\label{fig:failure_mode_pie}
\end{figure}

Ontology coverage provides a positive counterweight. The current ontology table shows 52 aspects tracked and 52 mapped to ontology paths. The highest-frequency mapped aspects are climate-detection with 79 records, governance with 66, missing with 60, and none with 23. Below these, domain-specific concepts such as roadmap karbon, pelatihan antikorupsi, komitmen net zero, and implementasi eco-mechanized mining are mapped to structured paths anchored to GRI-, governance-, or climate-oriented categories. This indicates that the ontology layer is currently stronger in coverage than the tone layer is in stability.

The denominator audit reinforces the diagnostic reading of these results. Of 5,444 extracted units in the relevant coverage view, 591 are excluded from tone distribution and agreement because they represent missing-tone cases. This is not a trivial cleanup issue. It shows that the extraction pipeline can remain operational while still failing on the central tone field often enough to distort downstream interpretation if those cases are ignored.

These findings reveal an important contradiction. The system is semantically organized enough to map diverse ESG aspects to ontology paths, yet it is still fragile when asked to consistently distinguish commitment, action, and outcome in noisy or governance-heavy contexts. That contradiction becomes one of the central interpretive findings of the chapter.

Overall, RQ3 is answered directly by the diagnostic evidence. The most important failure modes are missing tone, schema drift, hedged or modal language, and layout- or terminology-related ambiguity. OCR processing is operational at document level, but the central downstream weakness remains tone stability rather than ontology coverage.

\subsection{Quantitative Table Evidence as an Experimental Boundary}

One additional boundary must be stated explicitly for the Chapter 4 results: the present evidence layer is stronger for narrative extraction than for fully reconstructed quantitative tables. This does not mean the pipeline ignores tables; the failure-mode taxonomy already shows that table-heavy layouts contribute to extraction weakness. It does mean that KPI-bearing rows remain a sensitive error source whenever table structure is flattened, fragmented across pages, or detached from units and headers during preprocessing.

For that reason, the current chapter does not treat narrative extraction quality as a complete proxy for all ESG evidence quality. A stronger future benchmark should report table-specific metrics such as row reconstruction success, header continuity preservation, and metric-unit integrity on a reviewed subset. The revised thesis includes this boundary to avoid overstating the completeness of the current extraction evidence.

\subsection{RQ4: Stability and Reproducibility}

RQ4 asks to what extent ESG extraction outputs remain stable across varying prompts, LLM models, and service providers. This question is answered through prompt-level stability summaries, model-level summaries, stored workflow artifacts, and the reproducibility structure of the repository.

The repository performs strongly on reproducibility as an engineering system. The saved thesis dashboard report indexes 1,220 result artifacts and 184 LLM background jobs. Five static figures are already exported for thesis reuse: tone\_distribution.png, esg\_by\_tone.png, aspect\_by\_tone\_heatmap.png, climatebert\_label\_by\_tone.png, and climatebert\_remote\_top\_scores.png. In addition, the Chapter 4 Streamlit pages render live views over the same result tables, allowing the written thesis chapter and the interactive dashboard to remain aligned.

This reproducibility evidence matters because it shows that the research output is not dependent on one temporary notebook state. The file structure preserves both intermediate and summarized evidence. The pipeline can therefore be re-run, inspected, and extended without rebuilding the thesis reporting layer from scratch. At the same time, this should not be confused with a claim that all third-party LLM outputs are exactly repeatable across time. The thesis claims reproducible workflow structure and artifact persistence more strongly than deterministic output identity.

The strongest stability results appear in the model- and prompt-level summaries. Table~\ref{tab:model-level-extraction-performance} shows the current model comparison.

\begin{table}[htbp]
\centering
\caption{Model-level extraction performance.}
\label{tab:model-level-extraction-performance}
\footnotesize
\setlength{\tabcolsep}{4pt}
\renewcommand{\arraystretch}{1.15}
\begin{tabularx}{\linewidth}{|X|p{1.8cm}|X|p{1.8cm}|p{1.8cm}|p{1.8cm}|X|}
\hline
\textbf{Model} & \textbf{Runs} & \textbf{Parse success rate} & \textbf{Avg. records} & \textbf{Missing-tone rate} & \textbf{Schema-drift rate} & \textbf{Short interpretation} \\ \hline
\texttt{arcee-ai/trinity-large-preview:free} & 90 & 100.0\% & 3.02 & 0.05\% & 0.11\%  & Best stable thesis-facing baseline: high schema obedience and near-zero tone omission \\ \hline
\texttt{openai/gpt-oss-120b:free} & 20 & 100.0\% & 3.00 & 100.0\% & 42.5\%  & Formally parseable but practically unusable for tone analysis \\ \hline
\texttt{arcee-ai/trinity-large-thinking:free} & 89 & 89.9\% & 12.52 & 0.0\% & 0.0\%  & Very high extraction yield, but lower formal stability than the preview model \\ \hline
\texttt{minimax/minimax-m2.5:free} & 776 & 56.6\% & 4.94 & 0.05\% & 0.26\%  & High-volume live usage with weak parse reliability \\ \hline
\texttt{openai/gpt-oss-20b:free} & 145 & 95.9\% & 1.13 & 0.0\% & 0.0\%  & Structurally stable but low extraction yield \\ \hline
\texttt{qwen3.5:0.8b} & 2 & 0.0\% & 0.00 & 0.0\% & 0.0\%  & No usable extraction in the current slice \\ \hline
\end{tabularx}
\end{table}

The most important comparison for the thesis-facing stable subset is between arcee-ai/trinity-large-preview:free and openai/gpt-oss-120b:free. Both achieved perfect parse success in the summarized slice, but gpt-oss-120b exhibited a 100\% missing-tone rate and a 42.5\% schema-drift rate in the corresponding prompt-model grouping. This is a strong warning against treating model size or brand reputation as a proxy for task suitability. In this workflow, schema obedience and tone completion matter more than raw generative capability.

Figure~\ref{fig:model_tradeoff_scatter} complements Table~\ref{tab:model-level-extraction-performance} because it makes the trade-off between formal stability and usable extraction output easier to read.

\begin{figure}[H]
\begin{center}
  \includegraphics*[width=\textwidth]{model_tradeoff_scatter}
\end{center}
\caption{Model trade-off scatter plot with parse success on the x-axis and average extracted records on the y-axis. This figure helps separate models that are formally stable from those that are practically useful.}
\alt{Scatter plot comparing models by parse success and average extracted records, showing the trade-off between formal stability and practical extraction yield.}
\label{fig:model_tradeoff_scatter}
\end{figure}

Prompt-level stability shows a similar pattern. Chain-of-thought prompts, especially in English, produce the highest average records per run, while data.md produces formally parseable but functionally unusable tone outputs. Few-shot prompts appear inconsistent: the English few-shot prompt yields no average extracted records in the current summary, while the Indonesian few-shot prompt yields very low completion. This suggests that examples alone do not guarantee better extraction. In this repository, explicit tone framing and schema-constrained instruction style appear more important.

Figure~\ref{fig:prompt_strategy_comparison} should be read together with the prompt-level results table because prompt-family trends are central to the argument that prompt design matters more than nominal prompting style labels.

\begin{figure}[H]
\begin{center}
  \includegraphics*[width=\textwidth]{prompt_strategy_comparison}
\end{center}
\caption{Prompt-strategy comparison across average records, parse success, and field completion. The figure shows that chain-of-thought and tone-specific framing matter more than generic prompt validity alone.}
\alt{Three-panel comparison chart showing prompt templates across average extracted records, parse success, and field completion.}
\label{fig:prompt_strategy_comparison}
\end{figure}

These results reveal a practical trade-off. Some settings maximize extraction volume, while others maximize structural cleanliness. The best-performing configuration for this thesis is therefore not the one with the highest nominal complexity, but the one that balances parse success, tone completion, and manageable drift. This is exactly the type of engineering validity condition the chapter needs to report.

Overall, RQ4 is answered with an important qualification. The workflow is reproducible and stable in an engineering sense because prompts, outputs, logs, and derived artifacts are stored persistently, and because the same evaluation views can be regenerated. However, semantic output stability is not uniform across prompts, models, or providers. Stability depends strongly on prompt design and schema-following behavior, so the thesis can claim reproducible workflow structure more confidently than universally stable model behavior.

\section{Explainability Outputs}

The explainability layer helps show why the pipeline behaves as it does. At implementation level, code/rule\_based.py uses explicit lexical triggers for commitment, action, and outcome, while code/classical\_ml.py provides TF-IDF coefficient tables and local explanations. The lexicon file contains signals such as berkomitmen, commitment, will, menargetkan, target, aim to, and dedicated to for commitment-oriented language, and telah, achieved, has been, and successfully for outcome-oriented language.

The lexical trigger count summary confirms that these categories shape predictions in practice. Commitment triggers co-occurred 53 times with commitment predictions, compared with 30 action and 36 missing cases. Outcome triggers co-occurred 18 times with outcome predictions, but also 19 times with missing predictions and 17 times with commitment predictions. This helps explain why boundary cases remain difficult: many disclosure segments contain future-oriented and achieved language at the same time, especially in mixed narrative-reporting sections.

Two additional figures fit especially well in this subsection because they extend the explainability argument without requiring a new evaluation protocol.

\begin{figure}[htbp]
\begin{center}
  \includegraphics*[width=0.8\textwidth]{information_density_by_tone}
\end{center}
\caption{Information-density by tone. This figure shows whether commitment, action, outcome, and none records differ systematically in word count and therefore in extraction complexity.}
\alt{Boxplot comparing record word counts across commitment, action, outcome, and none tone categories.}
\label{fig:information_density_by_tone}
\end{figure}

\begin{figure}[H]
\begin{center}
  \includegraphics*[width=0.8\textwidth]{soft_language_ratio_by_tone}
\end{center}
\caption{Soft-language ratio by tone. This figure shows whether commitment records contain a higher proportion of soft or future-oriented verbs than action and outcome records, thereby giving a linguistic explanation for some tone-boundary failures.}
\alt{Boxplot comparing the ratio of soft or future-oriented verbs across commitment, action, and outcome tone categories.}
\label{fig:soft_language_ratio_by_tone}
\end{figure}

Ontology-path evidence is also interpretable. Environmental and climate-oriented examples include roadmap karbon, komitmen net zero, kerja sama energi bersih, and teknologi ramah lingkungan. Governance examples include pelatihan antikorupsi, kebijakan konflik kepentingan, and sertifikasi smap. These mappings show that the system can convert raw disclosure terms into structured thematic paths that are suitable for later chapter interpretation.

Representative records help illustrate the difference between commitment, action, and outcome:
\begin{itemize}
    \item Commitment example: “As a pioneer among industrial estate developers in Indonesia, BeFa is at the forefront of creating a green and eco-friendly industrial zone...” This record is labeled commitment, carries environmental claims, and receives a strong local climate-commitment prediction.
    \item Outcome example: “Keberhasilan pembangunan fasilitas perakitan kendaraan listrik di Magelang menjadi tonggak penting...” This record is labeled outcome with positive sentiment because it describes a realized infrastructure milestone
    \item Action example: “Pada tanggal 10 Oktober 2023, Perusahaan dan MB menandatangani perjanjian kerja sama...” This record is labeled action because it describes a concrete agreement, but not yet a completed performance outcome
\end{itemize}


These examples show that the tone taxonomy is meaningful in practice. The main challenge is not whether the categories are interpretable, but whether the extraction layer can assign them consistently across different prompts, models, and disclosure styles.

\section{Validation of Reference Data and Pilot Reference Layer}

The current evaluation target is not a full expert-labeled benchmark. It is a layered reference system built from extracted records, ClimateBERT-style comparison labels, weak lexical suggestions, and pilot human annotation files. This reference layer therefore has to be evaluated in its own right.

The strongest available comparison source is the ClimateBERT proxy agreement table over 332 records. It provides broad coverage and acceptable agreement, but it remains a proxy because the repository itself notes that a full one-to-one external ClimateBERT benchmark is still incomplete. The pilot human annotation layer is represented by files such as pilot\_ground\_truth\_seed.csv and pilot\_ground\_truth\_annotations.csv. The dashboard report currently records 70 pilot labels, which is useful for review and disagreement analysis but not enough to serve as a definitive gold standard.

Qualitatively, the reference layer is partially trustworthy but still noisy. Many pilot rows are marked needs\_review, which is appropriate and should be reported transparently. Several examples show disagreement between predicted tone and suggested tone. For instance, some governance or agreement-signing records are predicted as action but receive commitment-oriented heuristic suggestions because the text includes modal or forward-looking language. This is not merely annotation error; it reflects the conceptual difficulty of distinguishing action from commitment in governance-heavy sustainability language.

The repository also preserves invalid or weak reference cases rather than silently filtering them. Missing-tone and schema-drift rows remain visible in the review files. This is methodologically sound because it prevents the evaluation set from becoming artificially clean. At the same time, it means that the current reference layer should be described as a pilot reference framework rather than a finalized benchmark corpus.

Overall, the selected reference system is reliable enough for exploratory evaluation, disagreement analysis, and chapter-level empirical claims about patterns and failure modes. It is not yet reliable enough for definitive benchmark claims about absolute model superiority.

\section{Comparative Component Analysis}

The repository does not yet contain a full controlled ablation study in the strict machine-learning sense, where one component at a time is removed under a fixed benchmark and identical evaluation set. However, the existing results do support a comparative component analysis across prompts, models, and representational choices. This section therefore addresses the checklist requirement cautiously and explicitly.

The first comparative finding concerns the prompt layer. If prompt framing is treated as a removable component, then tone-specific chain-of-thought prompting is the most important component for usable T3 extraction. Replacing it with the generic data.md prompt preserves JSON parse success but causes catastrophic tone failure. This indicates that explicit tone instruction is more important than generic extraction framing for the thesis task.

The second finding concerns model choice. If model family is treated as a removable or swappable component, then not all large language models are equivalent. arcee-ai/trinity-large-preview:free produces stable parse and tone behavior in the thesis-facing subset, whereas openai/gpt-oss-120b:free shows unacceptable tone omission and schema drift in the summarized grouping. This means the pipeline depends more strongly on schema-following behavior than on raw model scale.

The third finding concerns representation type. The repository’s architecture suggests three feature regimes: handcrafted lexical features, sparse learned features, and contextual fused features. The rule-based model offers maximum interpretability and is useful for explainability and error diagnosis, but it is brittle under ambiguity. The classical TF-IDF model offers reproducible statistical baselines and coefficient-level explanations, but it is limited by sparse surface-form dependence. The hybrid contextual model is designed to capture sentence meaning, section context, and ontology information together, making it the most conceptually complete architecture. The current Chapter 4 evidence, however, is stronger for prompt/model stability than for full benchmark comparison among these three ABSA components, because the saved thesis dashboard focuses on pipeline evidence rather than controlled classifier benchmarking.

The fourth finding concerns ontology integration. The ontology layer appears robust in coverage, which implies that aspect-to-path mapping is not currently the most fragile component. If any component is most responsible for pipeline weakness, it is the tone-label production step rather than the ontology layer.

These component comparisons lead to a clear synthesis. The most important component of the implemented system is the prompt-and-schema design of the extraction layer, followed by the choice of model backend that can reliably honor that schema. Contextual and ontology-aware representations remain important for interpretability and future benchmarking, but the present evidence shows that extraction stability is the prerequisite for all later gains.

An additional integrative Sankey-style visual can be added later if the export is generated from the visualizer, because it would summarize how the most frequent aspects distribute across commitment, action, and outcome in one compact view.

\section{Chapter Synthesis}

The main findings of Chapter 4 are fivefold.

First, the pipeline successfully operationalizes PDF-to-structured ESG extraction over a nontrivial document set: 23 processed reports and roughly 5,512 OCR-covered pages. Second, the extraction layer produces a meaningful evidence store with 332 tone-bearing records and 2,074 T2 rows, confirming that the proposed schema is practically usable. Third, tone commitment aligns strongly with ClimateBERT-style commitment labels, with 83.7\% agreement and Cohen’s kappa of 0.645, which supports the construct relevance of the tone taxonomy without collapsing it into climate-topic classification. Fourth, the major weakness of the system is not parser failure but tone instability, especially missing-tone cases and schema drift under certain prompt settings. Fifth, the most important determinant of usable output is the interaction between prompt design and model compliance, not simply model scale.

The best overall thesis-facing configuration in the current saved evidence is the one represented by arcee-ai/trinity-large-preview:free with tone-focused prompts, because it combines high parse success with low missing-tone and low schema-drift rates. The most important system component is the tone-aware extraction prompt design. The main trade-off is between broad extraction flexibility and schema reliability: settings that appear generically capable do not necessarily produce thesis-usable structured outputs.

These findings connect directly back to Chapter 3. The methodology emphasized modularity, provenance, bilingual robustness, and interpretability. Chapter 4 shows that these design choices were justified. The system works best when it preserves traceable intermediate artifacts, separates tone from sentiment, and treats disagreement as information rather than noise. These results also prepare the next chapter: the discussion can now interpret commitment dominance, governance-related tone failures, model disagreement, and ontology coverage as substantive findings rather than mere implementation artifacts.






























































 
\chapter{Discussion}
  \label{discussion}
  



































\section{Key Findings and Synthesis}

\subsection{Main Findings and Their Meaning}

The central argument of this thesis is that ESG disclosure analysis for Indonesian sustainability reports becomes more informative when it is treated not only as a topic-detection problem, but as a tone-aware, provenance-preserving, ontology-aligned evidence extraction problem. The Chapter 4 results support this argument. The implemented workflow successfully transformed 23 sustainability-report PDFs into an auditable OCR-expanded corpus covering approximately 5,512 pages, generated 332 tone-bearing ESG records and 2,074 T2 rows, mapped 52 aspects to ontology paths, and achieved 83.7\% agreement with ClimateBERT-style commitment labels at Cohen’s kappa of 0.645.

The most important empirical finding is that disclosure tone is both measurable and meaningful. The extracted evidence set is dominated by commitment tone, with 115 of 332 records assigned to commitment, while environmental disclosure is the most frequent ESG pillar with 179 records. This means that the active corporate disclosure sample in the repository is not primarily organized around verified outcomes. It is organized around forward-looking commitments and strategic positioning. That pattern matters because it directly supports the thesis motivation: sentiment polarity alone would miss the difference between a positive achieved result and a positive promise.

At the same time, the thesis should state the contribution boundary more precisely. The tone layer is not presented here as a replacement for established disclosure-quality, disclosure-intention, or greenwashing frameworks. Its narrower contribution is to operationalize, at record level and with preserved provenance, whether a disclosure functions primarily as commitment, action, or realized outcome. In that sense, the framework refines existing interpretive approaches by adding a computationally explicit disclosure-stage distinction rather than by claiming a wholly separate theory of sustainability communication.

The strongest overall implementation configuration in the thesis-facing subset is the combination of tone-aware prompts and stable schema-following models, especially arcee-ai/trinity-large-preview:free with tone-specific prompt templates. This configuration does not simply produce more text. It produces structurally usable records with very low missing-tone and schema-drift rates relative to weaker alternatives. By contrast, openai/gpt-oss-120b:free achieved perfect parse success in the summarized slice but produced a 100\% missing-tone rate and a 42.5\% schema-drift rate in the corresponding grouping. The meaning of this result is clear: for this task, schema compliance and tone completion are more important than nominal model scale or general fluency.

The most important technical insight from Chapter 4 is therefore not that a specific model “won,” but that the decisive factor is the interaction between prompt design and model behavior. The thesis methodology assumed that prompt sensitivity and structured-output fragility would be major determinants of quality. The Chapter 4 results confirm that assumption. They also show that ontology coverage is comparatively robust, while tone assignment remains the most fragile part of the system. This partly confirms and partly sharpens Chapter 3. Chapter 3 anticipated modular weakness; Chapter 4 demonstrates that the main bottleneck is concentrated in tone stability rather than in OCR completion or ontology mapping alone.

\begin{figure}[ht]
\begin{center}
  \includegraphics*[width=\textwidth]{tone_distribution}
\end{center}

\caption{Tone distribution in the active thesis-facing subset. This figure is used to interpret commitment dominance as a substantive disclosure pattern rather than as a descriptive count reported in Chapter 4.}\alt{Bar chart showing the tone distribution of extracted records in the thesis-facing subset, with commitment as the largest category.}
\label{fig:discussion_tone_distribution}
\end{figure}

\subsection{Deeper Patterns and Research Question Mapping}

The deeper pattern behind the results is that different components of the system rely on different kinds of information. The OCR layer depends on document completeness and page-level traceability. The extraction layer depends on prompt framing, model obedience, and output schema stability. The comparison layer depends on construct overlap between tone and climate-topic signals. The ontology layer depends on semantic coverage of ESG vocabulary. These components do not fail for the same reason, and that is why the modular design matters.

The most influential component in the full pipeline is the tone-aware extraction prompt design. This can be seen directly in the prompt stability table: data.md is parseable but functionally poor for tone, whereas chain-of-thought tone prompts yield far better record counts and lower omission. This means that the pipeline’s best results are not produced by generic extraction capacity, but by explicit instruction that forces the model to separate commitment, action, and outcome.

Different methods also rely on different information types. The rule-based model relies on surface lexical cues such as berkomitmen, will, target, achieved, and successfully. The classical TF-IDF model relies on sparse word and character statistics. The hybrid model is designed to rely on contextual and ontology-aware information. The LLM extraction layer, however, relies on broader contextual reasoning and instruction following. This explains why some governance-heavy or agreement-signing texts become difficult: they contain both action markers and commitment framing, so simple lexical cues and broad generative interpretation can disagree.

The following matrix summarizes how the main research questions map to evidence.

All four research questions are therefore answered in the current thesis scope, although not all answers are equally mature. RQ1 and RQ4 are strongest at operational and diagnostic level because the repository preserves provenance, artifacts, and visible failure cases rather than hiding them. RQ2 and RQ3 are also supported, but they still depend on pilot-level validation and proxy comparison rather than full gold-standard benchmarking.

\begin{table}[htbp]
\centering
\caption{Research-question evidence summary.}
\label{tab:research-question-evidence-summary}
\footnotesize
\setlength{\tabcolsep}{4pt}
\renewcommand{\arraystretch}{1.15}
\begin{tabularx}{\linewidth}{|p{2cm}|X|p{2cm}|p{2cm}|X|}
\hline
\textbf{Research question} & \textbf{Main Chapter 4 evidence} & \textbf{Figure or Table} & \textbf{Current confidence level} & \textbf{Remaining limitation} \\ \hline
\texttt{RQ1. ESG ABSA Schema} & 23 processed OCR documents, about 5,512 pages, 332 tone-bearing records, 2,074 T2 rows, and 52 mapped ontology aspects & Figure~\ref{fig:discussion_tone_distribution}, Figure~\ref{fig:esg_by_tone}, Figure~\ref{fig:aspect_by_tone_heatmap} & Moderate to high & Evidence is operationally strong, but class balance remains uneven and social coverage is weak \\ \hline
\texttt{RQ2. Tone vs. Climate-Specific Models} & 83.7\% agreement and Cohen's kappa of 0.645 between tone commitment and climate-commitment proxy labels across 332 records & Table~\ref{tab:tone-commitment-versus-climatebert-style-commitment-comparison} and Figure~\ref{fig:discussion_climatebert_by_tone} & Moderate & ClimateBERT comparison is still proxy-based rather than a complete one-to-one external benchmark \\ \hline
\texttt{RQ3. Pipeline Diagnostics} & 61 missing-tone failures, 20 schema-drift cases, and a concentrated failure distribution dominated by a small set of recurring weaknesses & Table~\ref{tab:failure-mode-count-table}, Figure~\ref{fig:failure_mode_pareto}, Figure~\ref{fig:failure_mode_pie} & High for diagnostic claims & Failure analysis is strong, but some categories still need more manually curated case examples \\ \hline
\texttt{RQ4. Stability and Reproducibility} & 1,220 stored result artifacts, 184 background jobs, prompt-level summaries, and clear model-level trade-offs between parse stability and usable extraction & Table~\ref{tab:prompt-level-extraction-performance-across-the-thesis-facing-subset}, Table~\ref{tab:model-level-extraction-performance}, Figure~\ref{fig:model_tradeoff_scatter}, Figure~\ref{fig:prompt_strategy_comparison} & Moderate to high & Workflow reproducibility is strong, but semantic outputs remain sensitive to prompt/model/provider changes \\ \hline
\end{tabularx}
\end{table}

\section{Research Question Resolution}

\subsection{Resolution of RQ1 and RQ2}

RQ1 asks whether sustainability-report PDFs can be transformed into structured ESG evidence through the implemented pipeline. The answer is yes. The OCR and preprocessing workflow processed 23 tracked documents and approximately 5,512 pages, and the downstream extraction pipeline produced a stable evidence store with 332 tone-bearing records and 2,074 T2 outputs. In practical terms, this means the thesis has demonstrated the feasibility of moving from raw report files to page-linked, machine-readable ESG evidence at nontrivial corpus scale.

RQ2 asks whether the system can represent ESG disclosure beyond coarse topic detection by separating pillar, aspect, sentiment, and disclosure tone. The answer is also yes, but with an important qualification. The active result set shows meaningful distributions across tone and ESG pillar, with commitment as the dominant tone and environmental disclosure as the dominant pillar. This confirms that the schema can capture multiple dimensions of disclosure behavior. However, the pilot annotation evidence and the missing-tone cases show that this representation is not yet equally reliable across all text types. The contribution of RQ2 is therefore not only that the schema exists, but that the thesis demonstrates where it is robust and where it needs stronger human calibration.

\subsection{Resolution of RQ3 and RQ4}

RQ3 asks how the tone-aware pipeline compares with ClimateBERT-style climate classification. The answer is that the two are strongly related but not interchangeable. The agreement table shows 83.7\% percent agreement and Cohen’s kappa of 0.645 between tone commitment and climate-commitment labels across 332 records. This is a strong descriptive result because it validates the relevance of the tone taxonomy while also preserving its distinctiveness. In practical terms, ClimateBERT-style labels can tell us that a disclosure is climate-related, but tone labels can additionally tell us whether the disclosure reports a promise, an action, or a realized outcome.

\begin{figure}[ht]
\begin{center}
  \includegraphics*[width=\textwidth]{climatebert_label_by_tone}
\end{center}
\caption{Relationship between LLM tone labels and ClimateBERT-style labels. Reused here because the discussion interprets this cross-distribution as evidence of partial construct overlap rather than full label equivalence.}
\alt{Chart showing the relationship between tone categories and ClimateBERT-style labels to illustrate partial construct overlap.}
\label{fig:discussion_climatebert_by_tone}
\end{figure}

RQ4 asks what weaknesses remain in the current evaluation and extraction strategy. The answer is that the major weaknesses are concentrated in tone omission, schema drift, and ambiguity-rich text. The strongest evidence is the 61 missing-tone cases, the 19 schema-drift cases in missing predictions, the hedged-or-modal-language failures, and the sensitivity of results to prompt and model choice. The proposed evaluation strategy is sufficient to reveal these weaknesses because it combines OCR completion, parse success, stability summaries, disagreement tables, ontology coverage, and review-oriented pilot files. However, it is not yet sufficient to claim final benchmark superiority because the gold-standard annotation layer remains limited. Taken together, the four RQs show that the thesis objective has been achieved at the level of an executable, interpretable, and diagnostically rich ESG disclosure-analysis framework.



\section{Limitations}

The first major limitation is the reference-data limitation. The thesis does not yet rely on a complete expert-labeled benchmark. Instead, it uses a layered reference system built from extracted records, ClimateBERT-style comparison labels, heuristic suggestions, and pilot human annotation files. This is sufficient for exploratory evaluation and disagreement analysis, but it limits the strength of claims about absolute model performance.

The second limitation is domain and sample concentration. The results are grounded in Indonesian sustainability and annual reports, with bilingual or mixed-language disclosure patterns. This is a strength for domain specificity, but it limits immediate generalization to other regulatory settings, other document genres, or monolingual corpora. The social pillar is also underrepresented in the current thesis-facing subset, which means tone and aspect findings are shaped more strongly by environmental and governance disclosure.

The third limitation is model and prompt sensitivity. The prompt stability table shows that formally valid prompts can still produce unusable tone outputs. The model stability table shows that some models with perfect parse success can still fail semantically by omitting tone or drifting schema. This means that the system is not yet robust to arbitrary backend substitution. Its performance depends on careful pairing between prompt design and model behavior.

The fourth limitation is the existence of technical failure cases that affect interpretation. The failure-mode table includes 61 missing-tone outputs, 19 schema-drift cases, hedged-language failures, table-layout problems, and issues with regulatory Indonesian terms. These failures do not invalidate the pipeline, but they do mean that the Chapter 4 findings should be interpreted as evidence of a strong but still imperfect research system.

The fifth limitation is that the greenwashing index remains heuristic. Company-level ratios such as 73.0 for BEST-SR\_2024 or 4.778 for vktr\_ar\_sr\_2024 are analytically interesting, but they are not yet validated against external ESG ratings, enforcement cases, or human expert judgment. For that reason, the index should currently be interpreted as a screening or discussion aid, not as a final adjudication tool.

Two figures fit well in this limitations section because they help the reader see where the evidence is uneven rather than only being told so.

\begin{figure}[ht]
\begin{center}
  \includegraphics*[width=\textwidth]{greenwashing_gap_scatter}
\end{center}
\caption{Commitment-outcome gap by company. This figure visualizes the relationship between commitment count, outcome count, and the company-level screening ratio, while the surrounding text keeps clear that the score remains heuristic.}
\alt{Scatter plot showing company-level commitment counts versus outcome counts, with point size or color indicating the screening ratio.}
\label{fig:greenwashing_gap_scatter}
\end{figure}

\begin{figure}[ht]
\begin{center}
  \includegraphics*[width=\textwidth]{commitment_outcome_ratio}
\end{center}
\caption{Tone-share ratio chart for the active extracted-record layer. This figure supports the discussion claim that commitment-heavy evidence remains much more common than outcome-heavy evidence in the current subset.}
\alt{Bar chart showing the share of extracted records across tone categories, highlighting the higher share of commitment records relative to outcome records.}
\label{fig:commitment_outcome_ratio}
\end{figure}


\section{Future Work}
\subsection{Immediate Technical Extensions}

The most immediate improvement is to expand and formalize the human annotation layer. The current pilot annotation files should be extended into a stratified expert-reviewed benchmark with inter-annotator agreement, especially for commitment versus action boundary cases, governance-heavy disclosures, and records currently marked needs\_review. This directly addresses the weakest part of the evaluation design.

The second immediate improvement is to strengthen tone-specific prompt engineering and schema enforcement. The results already show that tone-specific chain-of-thought prompts outperform generic extraction prompts. The next step is to convert this finding into a more systematic prompt protocol, including stricter output validation, field-level confidence checks, and targeted rerun logic for missing-tone outputs.

The third immediate improvement is OCR quality evaluation. The repository currently proves that OCR completion is operational, but it does not yet report character error rate, table extraction accuracy, or page-level semantic quality baselines. Adding those checks would help separate upstream OCR noise from downstream tone-classification difficulty.

The fourth immediate improvement is one-to-one ClimateBERT benchmarking over the full extracted record set. The repository already supports record-batch generation for this purpose. Completing that pipeline would turn the current proxy comparison into a stronger external reference layer.

The fifth immediate improvement is better filtering of non-ESG and accounting-policy text. The pilot annotation files show examples where financial or accounting statements are correctly treated as none, but more explicit front-end filtering could reduce wasted extraction capacity and improve downstream quality.

\subsection{Long-Term Research Roadmap}

In the longer term, this research can grow into a broader program on computational sustainability disclosure analysis for low-resource and bilingual settings. One direction is model training: the hybrid contextual architecture could be extended through supervised fine-tuning on a larger Indonesian ESG tone corpus, allowing the system to learn tone distinctions more directly rather than depending so heavily on prompt-driven extraction.

A second direction is robustness and domain transfer. The framework should be tested on longer, noisier, and more sectorally diverse reports, and then extended to other regulatory environments beyond the current Indonesian focus. This would reveal whether the tone taxonomy and ontology structure remain useful under different reporting cultures.

A third direction is deployment-oriented research. Because the repository already stores page-level provenance and rich artifacts, it could be adapted into an analyst-facing review system for regulators, ESG research teams, investors, or auditors. In that setting, the most important research question would shift from “can the model classify?” to “how should humans and models collaborate when evidence quality is uncertain?”

A fourth direction is graph and knowledge-integration research. The ontology and semantic-export components already point toward graph-based querying and GraphRAG-style workflows. Future work could examine how structured ESG evidence, ontology paths, and source-page provenance can support traceable retrieval rather than only static reporting.

Taken together, these future directions form a coherent roadmap: first stabilize annotation and benchmarking, then improve model robustness and OCR quality, then extend across domains and deployment contexts, and finally integrate the workflow into broader evidence-centric research infrastructure.

\section{Broader Impact}
\subsection{Practical and Societal Benefits}

The most direct beneficiaries of this research are ESG analysts, sustainability-report authors, regulators, investors, and research teams working on corporate accountability. For ESG analysts, the system reduces the cost of moving from hundreds of PDF pages to a smaller set of structured, reviewable disclosure records. For sustainability-report authors, the tone framework provides feedback about whether a report is dominated by commitments, operational actions, or realized outcomes. For regulators and oversight bodies, the provenance-preserving design offers a way to inspect disclosure claims without losing the link to the source page. For investors, the greenwashing-oriented ratios and tone distributions can serve as a supplementary signal when reading large volumes of narrative disclosure.

The broader practical value of the thesis is that it improves on existing ESG automation in three ways. First, it preserves auditable page-level lineage. Second, it distinguishes disclosure stage rather than treating all positive ESG language as equivalent. Third, it reveals failure modes and uncertainty instead of hiding them behind a single summary score. These are meaningful improvements for real-world use because ESG decision-making often depends as much on traceability and interpretability as on aggregate classification accuracy.

At societal level, the work supports stronger information transparency in sustainability communication. If deployed responsibly, a system like this can help identify when corporations communicate mostly aspiration, when they report concrete action, and when they document measurable outcomes. That distinction is valuable for accountability, sustainability governance, and public understanding of corporate climate and ESG claims.

\subsection{Ethical Risks, Bias, and Responsible Use}

The main ethical risk is over-automation. If users treat the extracted tone labels or greenwashing ratios as definitive truth rather than as decision-support signals, the system could mischaracterize organizations or obscure nuance in legally or strategically sensitive disclosures. This risk is especially important because the current evaluation layer is still partly weakly supervised.

Bias and fairness concerns also remain. The current active evidence layer overrepresents environmental and governance language and underrepresents the social pillar. The model and prompt outputs are also sensitive to bilingual phrasing, governance boilerplate, and table-heavy layouts. Without careful review, the system may systematically underperform on certain sectors, disclosure styles, or language patterns.

Privacy concerns are relatively limited because the corpus is based on public corporate reports, but responsible deployment still requires care around data handling, caching, provenance logs, and model outputs. Misuse is more likely to arise through misinterpretation than through private-data exposure. For example, a high greenwashing index should not be treated as proof of misconduct without human review and external corroboration.

The appropriate safeguards are therefore clear: maintain source-page traceability, preserve uncertainty and failure flags, require human review for sensitive cases, expand annotation before high-stakes deployment, and avoid using the system as a sole basis for regulatory or investment judgment. In sustainability terms, the work is aligned with broader responsible-governance goals because it encourages more transparent, evidence-based reading of ESG claims. The balanced conclusion is that this is useful technology for ESG evidence organization and interpretation, but its value depends on disciplined, review-oriented, and context-aware use.

An additional deployment-oriented chart can still be added later if you export a reliability-adjusted waterfall or a qualitative case-study figure for one representative company. If used, the surrounding text must make clear that the output is illustrative and not a finalized evaluative instrument.
   
\chapter{Conclusion}
  \label{summary}
  This thesis set out to design an executable framework for analysing Indonesian sustainability reports as structured ESG evidence rather than as unstructured narrative text. The core objective was not only to classify disclosures, but to build a reproducible workflow that transforms heterogeneous PDF reports into auditable records with page-level provenance, aspect labels, ESG pillar assignments, sentiment, and disclosure tone. This objective was achieved through a modular pipeline that combined OCR, page-aware extraction, LLM-based structuring, ontology alignment, comparison against ClimateBERT-style labels, and dashboard-based diagnostics.

The results show that the proposed approach is technically feasible and analytically useful. The active thesis subset covered 23 processed reports and about 5,512 pages, producing 332 tone-bearing ESG records and 2,074 downstream rows. The record-level schema was operational, all 52 tracked aspects were mapped to ontology paths, and the comparison with ClimateBERT-style commitment labels showed strong but not perfect alignment, with 83.7

A key methodological choice in this thesis was to treat the system as an executable research workspace rather than as a single-model benchmark. This choice was justified by the results. The main strengths of the framework are traceability, modularity, and diagnostic transparency. The main weaknesses are also clear: tone extraction remains sensitive to prompt design and model behavior, and missing-tone cases and schema drift show that fully reliable automation has not yet been reached. These limitations do not reduce the value of the work; instead, they show where future validation and benchmark development are needed.

The significance of the thesis is practical as well as academic. For researchers, it provides a reproducible basis for ESG disclosure analytics in a bilingual and low-resource setting. For analysts, regulators, and reporting teams, it offers a way to move from long PDF reports to structured, reviewable evidence without losing the connection to the source pages. More broadly, the work shows that useful ESG text analytics requires not only better models, but also better document engineering, provenance preservation, and evaluation design. In that sense, the thesis contributes an engineering-oriented foundation for more transparent and auditable sustainability-report analysis.     
\printbibliography[heading=bibnumbered, title=\mybibname]

\startappendix
\chapter{Appendices}
  































This appendix chapter consolidates supplementary material that supports the methodological and reproducibility claims in the main thesis. The appendices are intended to document materials that are important for transparency but too detailed to place in the core chapter narrative.

\section{Appendix Contents}\label{appendix-contents}

The appendices should contain, at minimum, the following materials:

\begin{enumerate}
\def\labelenumi{\arabic{enumi}.}
\tightlist
\item
  prompt templates used in the main thesis-facing comparisons;
\item
  model identifiers, providers, and decoding settings where available;
\item
  benchmark and comparison artifact paths;
\item
  annotation instructions and tone-label definitions;
\item
  additional difficult examples and disagreement cases;
\item
  supplementary failure-mode tables;
\item
  extended reproducibility notes for rerunning stored analyses.
\end{enumerate}

\section{Tone-Label Definitions}\label{tone-label-definitions}

The pilot annotation and review process should use the following working definitions:

\begin{enumerate}
\def\labelenumi{\arabic{enumi}.}
\tightlist
\item
  \textbf{Commitment:} forward-looking promises, targets, intentions, or policy declarations that do not yet demonstrate completed implementation or realized outcome.
\item
  \textbf{Action:} concrete activities, agreements, programs, investments, or operational steps that are being undertaken or have been initiated.
\item
  \textbf{Outcome:} realized, completed, or demonstrably achieved results, milestones, or measured performance changes.
\item
  \textbf{Needs review:} records whose language is too ambiguous, mixed, or incomplete for stable tone assignment without additional human interpretation.
\end{enumerate}

\section{Reproducibility Notes}\label{reproducibility-notes}

The main reproducibility artifacts referenced in the thesis include:

\begin{enumerate}
\def\labelenumi{\arabic{enumi}.}
\tightlist
\item
  OCR-expanded document folders under \texttt{data/thesis\_dataset/};
\item
  extracted ESG records under \texttt{results/esg\_records.json};
\item
  resumable benchmark outputs under \texttt{results/t1\_results.jsonl} and \texttt{results/t2\_results.jsonl};
\item
  prompt and model stability summaries under \texttt{results/revision\_analysis/};
\item
  dashboard-ready figures and tables under \texttt{results/thesis\_workflow\_dashboard/}.
\end{enumerate}

These artifacts support rerunning, auditing, and extending the workflow even where exact third-party LLM outputs may vary across time.

\section{End-to-End User Workflow Diagram}\label{end-to-end-user-workflow-diagram}

The appendix should also document the operational user flow across the main thesis-facing tools. In the implemented repository, the workflow begins with OCR preparation in \texttt{pages/Bulk\_OCR.py}, continues through page-range-aware ESG extraction in \texttt{pages/llm\_processing.py}, and then passes through the ClimateBERT comparison stage as the T1 pipeline.

\subsection{Workflow Summary}\label{workflow-summary}

The implemented user flow is:

\begin{enumerate}
\def\labelenumi{\arabic{enumi}.}
\tightlist
\item
  upload or select PDF files in \textbf{Bulk OCR}
\item
  run OCR and save page-level outputs into \texttt{data/thesis\_dataset/}
\item
  open \textbf{LLM Processing}
\item
  select the OCR-expanded PDF/document target
\item
  choose a page range for processing
\item
  choose one of three LLM-provider families:

  \begin{itemize}
  \tightlist
  \item
    \textbf{OpenRouter}
  \item
    \textbf{LM Studio / OpenAI-compatible}
  \item
    \textbf{Ollama}
  \end{itemize}
\item
  run structured ESG extraction over the selected page range
\item
  save structured records into \texttt{results/esg\_records.json}
\item
  run \textbf{ClimateBERT} processing as the T1 comparison layer
\item
  save ClimateBERT outputs into \texttt{results/predictions.json} or resumable T1 artifacts
\end{enumerate}


\begin{figure}[H]
\begin{center}
  \includegraphics*[width=\textwidth]{mermaid_flow_diagram_01}
\end{center}
\caption{Bulk OCR to LLM-processing flow diagram. This figure visualizes the operational path from file intake, OCR expansion, and page storage to provider-specific ESG extraction and downstream ClimateBERT comparison outputs.}
\alt{Flowchart showing user entry through Bulk OCR, OCR output storage, document selection in LLM Processing, provider choice among OpenRouter, LM Studio, and Ollama, ESG record storage, and downstream ClimateBERT output generation.}
\label{fig:appendix_bulkocr_llm_climatebert_flow}
\end{figure}

\subsection{Mermaid Sequence Diagram}\label{mermaid-sequence-diagram}

\begin{figure}[H]
\begin{center}
  \includegraphics*[width=\textwidth]{mermaid_flow_diagram_02}
\end{center}
\caption{Bulk OCR to ClimateBERT sequence diagram. This figure shows the interaction order between the user, OCR storage, LLM processing, provider backends, ESG record storage, and the downstream ClimateBERT comparison stage.}
\alt{Sequence diagram showing the user uploading or selecting files, OCR artifact creation, page-range selection in LLM Processing, requests to an LLM provider, ESG record persistence, and ClimateBERT prediction storage.}
\label{fig:appendix_bulkocr_llm_climatebert_sequence}
\end{figure} 

Sequence explanation:

\begin{enumerate}
\def\labelenumi{\arabic{enumi}.}
\tightlist
\item
  The user starts in \textbf{Bulk OCR}, either by uploading files through the browser or by selecting server-side PDFs already copied into the working directory.
\item
  The OCR page sends the selected files to the OCR service and stores the returned artifacts as document folders containing \texttt{ocr\_result.json}, page markdown files, and image crops.
\item
  The user then moves to \textbf{LLM Processing}, which loads one OCR-expanded document at a time.
\item
  Inside \textbf{LLM Processing}, the user selects a page range rather than always processing the full report. This is important because the extraction workflow is designed around page batches.
\item
  The user chooses one provider family for the extraction stage: \textbf{OpenRouter}, \textbf{LM Studio / OpenAI-compatible}, or \textbf{Ollama}.
\item
  The selected page batch is sent to the chosen LLM provider, and the returned output is parsed into the ESG record schema.
\item
  Structured records are appended into \texttt{results/esg\_records.json}, which becomes the main evidence layer for later analysis.
\item
  The extracted text layer is then passed into the \textbf{ClimateBERT} T1 stage, which acts as a downstream comparison or benchmark layer rather than as the first ingestion step.
\item
  ClimateBERT-style predictions are saved into \texttt{results/predictions.json} or resumable T1 artifacts such as \texttt{results/t1\_results.jsonl}.
\end{enumerate}

\subsection{LaTeX Figure Version}\label{latex-figure-version}

The same operational workflow can also be represented in a LaTeX-ready figure block for the final thesis.

\begin{figure}[H]
\begin{center}
  \includegraphics*[width=\textwidth]{mermaid_flow_diagram_01}
\end{center}
\caption{Appendix workflow from Bulk OCR to page-range-based LLM extraction and downstream ClimateBERT processing. The implemented flow begins with OCR expansion, continues with provider-specific ESG extraction over selected page ranges, and ends with ClimateBERT comparison over the extracted text layer.}
\alt{Workflow diagram showing user upload or file selection in Bulk OCR, OCR storage into thesis dataset folders, page-range selection in LLM Processing, choice among OpenRouter, LM Studio, or Ollama, storage of structured ESG records, and downstream ClimateBERT T1 processing.}
\label{fig:appendix_bulkocr_llm_climatebert_workflow}
\end{figure}

\subsection{Interpretation}\label{interpretation}

This diagram shows that ClimateBERT processing is not the first-stage ingestion tool. Instead, it is a downstream comparison or benchmark layer applied after OCR expansion and after page-range-based ESG extraction. That distinction is methodologically important because the thesis does not compare ClimateBERT directly against raw PDFs. It compares ClimateBERT-style predictions against the extracted text layer produced by the broader ESG pipeline.

\subsection{Tone-Prompt Family Visualization and Prompt Reasoning}\label{tone-prompt-family-visualization-and-prompt-reasoning}

The repository contains six thesis-facing tone prompts under \texttt{prompt/}:

\begin{itemize}
\tightlist
\item
  \texttt{tone\_zero\_shot\_english.md}
\item
  \texttt{tone\_zero\_shot\_indonesian.md}
\item
  \texttt{tone\_few\_shot\_english.md}
\item
  \texttt{tone\_few\_shot\_indonesian.md}
\item
  \texttt{tone\_chain\_of\_thought\_english.md}
\item
  \texttt{tone\_chain\_of\_thought\_indonesian.md}
\end{itemize}

These prompt files share the same extraction target, but they differ in how much reasoning structure they impose before the final JSON output. The prompt family is summarized below so that the appendix makes clear that prompt variation is not cosmetic. It is one of the main experimental levers of the thesis.

\subsubsection{Prompt-Family Logic}\label{prompt-family-logic}



\subsubsection{Prompt-Family Summary Table}\label{prompt-family-summary-table}

\begin{table}[htbp]
\centering
\caption{Tone-prompt family summary.}
\label{tab:tone-prompt-family-summary}
\footnotesize
\setlength{\tabcolsep}{4pt}
\renewcommand{\arraystretch}{1.15}
\begin{tabularx}{\linewidth}{|p{0.18\linewidth}|p{0.14\linewidth}|X|X|X|}
\hline
\textbf{Prompt file} & \textbf{Prompt style} & \textbf{Main reasoning structure} & \textbf{Intended advantage} & \textbf{Main risk} \\ \hline
\texttt{tone\_zero\_shot\_english.md} & Zero-shot, English & Direct instructions, label list, tone priority rules, JSON schema & Strong baseline for testing whether explicit definitions alone are enough & May remain too shallow for ambiguous governance or mixed-tone text \\
\texttt{tone\_zero\_shot\_indonesian.md} & Zero-shot, Indonesian & Same logic as English version, but adapted to Indonesian cue words and examples of commitment/action/outcome distinction & Better alignment with Indonesian report language and bilingual segments & Longer instruction set can still be ignored when outputs drift \\
\texttt{tone\_few\_shot\_english.md} & Few-shot, English & Adds worked examples for commitment, outcome, and risk-oriented none tone & Demonstration-based anchoring for label and tone boundaries & Examples may overfit narrow patterns and suppress extraction breadth \\
\texttt{tone\_few\_shot\_indonesian.md} & Few-shot, Indonesian & Adds Indonesian worked examples and local wording for target, risk, and measured result & Better grounding for Indonesian narrative style and common lexical cues & Example dependence may still fail on governance-heavy or non-example structures \\
\texttt{tone\_chain\_of\_thought\_english.md} & Chain-of-thought style, English & Hidden internal steps for segmentation, tone assignment, sentiment assignment, and final schema fill & Forces the model to separate reasoning stages before output & Longer prompt can increase verbosity pressure or schema drift in weaker models \\
\texttt{tone\_chain\_of\_thought\_indonesian.md} & Chain-of-thought style, Indonesian & Most explicit reasoning scaffold, local cue words, tone priority logic, and concise reasoning guidance & Best fit for difficult Indonesian and mixed-language extraction conditions & Most complex prompt, so weaker backends may fail to follow every instruction consistently \\
\hline
\end{tabularx}
\end{table}

\subsubsection{Why the Prompt Reasoning Matters}\label{why-the-prompt-reasoning-matters}

The central prompt-design idea is not just to ask for JSON, but to force the model to reason about disclosure function before assigning sentiment. In this thesis, that means the prompt must first distinguish:

\begin{enumerate}
\def\labelenumi{\arabic{enumi}.}
\tightlist
\item
  whether a segment is explicitly ESG-relevant;
\item
  what aspect and label family it belongs to;
\item
  whether the statement is a commitment, action, or outcome;
\item
  only after that, whether the statement carries positive, negative, neutral, or none sentiment.
\end{enumerate}

This ordering matters because many sustainability disclosures are future-oriented promises written in positive language. If the prompt does not explicitly disentangle tone from sentiment, the model can collapse a promise into a falsely positive achieved result. The tone prompts therefore use reasoning scaffolds to reduce this specific validity error.

\subsubsection{Prompt Reasoning Description by Family}\label{prompt-reasoning-description-by-family}

\textbf{Zero-shot reasoning.} The zero-shot prompts rely on explicit definitions, priority rules, and schema constraints without demonstration examples. Their reasoning philosophy is that tone extraction should remain possible from a strong instruction set alone. In methodological terms, zero-shot prompts act as the cleanest test of whether the schema itself is intelligible to the model.

\textbf{Few-shot reasoning.} The few-shot prompts add short worked examples so the model sees what commitment, outcome, and climate-risk cases should look like in JSON form. Their reasoning philosophy is analogical: the model is guided to map new inputs to demonstrated output patterns. This is useful for ambiguous ESG phrasing, but it can also narrow the extraction behavior if the examples are too stereotyped.

\textbf{Chain-of-thought-style reasoning.} The chain-of-thought prompts explicitly instruct the model to segment text, identify ESG relevance, assign tone, then assign sentiment, while keeping the actual reasoning hidden from the final response. Their reasoning philosophy is procedural: a difficult extraction problem is decomposed into smaller internal steps so that tone and sentiment are less likely to be conflated. This is the most thesis-aligned prompt family because the research question depends on stable tone disentanglement rather than generic JSON completion alone.

\subsubsection{Appendix Interpretation}\label{appendix-interpretation}

This prompt family should be read as a controlled prompt-engineering ladder rather than as six unrelated files. The zero-shot prompts test direct schema interpretability. The few-shot prompts test example anchoring. The chain-of-thought prompts test whether stepwise internal reasoning improves tone stability. That is why prompt variation in Chapter 4 is analytically meaningful: it probes whether the thesis contribution depends only on model choice or also on reasoning structure embedded in the prompt itself.

\section{JSON Data Examples}\label{json-data-examples}

The repository uses JSON extensively for OCR outputs, extraction records, dashboard artifacts, model caches, logs, and workflow decisions. This appendix section provides representative examples of the actual JSON structures used by the implemented system.

\subsection{OCR Output JSON}\label{ocr-output-json}

Path example:

\begin{itemize}
\tightlist
\item
  \texttt{data/thesis\_dataset/2017\ Sustainability\ Report\ PT\ Bank\ Permata\ Tbk\ (Rev\ May2018)\_0\_pdf/ocr\_result.json}
\end{itemize}

Purpose:

\begin{itemize}
\tightlist
\item
  stores page-level OCR output
\item
  preserves extracted markdown text
\item
  stores image coordinates and embedded image payloads
\item
  records OCR model and usage metadata
\end{itemize}

Example:

\begin{verbatim}
{
  "pages": [
    {
      "index": 0,
      "markdown": "PermataBank\n\nLaporan Keberlanjutan 2017...",
      "images": [
        {
          "id": "img-0.jpeg",
          "top_left_x": 48,
          "top_left_y": 260,
          "bottom_right_x": 667,
          "bottom_right_y": 777,
          "image_base64": "data:image/jpeg;base64,..."
        }
      ]
    }
  ],
  "model": "...",
  "document_annotation": "...",
  "usage_info": "..."
}
\end{verbatim}

\subsection{ESG Extraction Run JSON}\label{esg-extraction-run-json}

Path example:

\begin{itemize}
\tightlist
\item
  \texttt{results/esg\_records.json}
\end{itemize}

Purpose:

\begin{itemize}
\tightlist
\item
  stores T3 LLM extraction runs
\item
  preserves prompt, model, page target, records, and failure state
\item
  acts as the main structured ESG evidence store
\end{itemize}

Success example:

\begin{verbatim}
{
  "timestamp": "2026-06-03T18:38:40Z",
  "model": "nvidia/nemotron-3-nano-30b-a3b:free",
  "target": "Sustainability_Report_PT_Sentul_City_Tbk_2025_pdf/batch_7",
  "target_pages": "page_0012.md, page_0013.md",
  "prompt": "tone_chain_of_thought_indonesian.md",
  "ok": true,
  "records": [
    {
      "text": "Elevating Lives, Building Sustainable Growth",
      "aspect": "strategic vision",
      "labels": ["climate-d", "climate-commitment"],
      "esg": "E",
      "tone": "commitment",
      "sentiment": "positive",
      "sentiment_score": 1,
      "reasoning": "The phrase expresses a forward-looking ..."
    }
  ]
}
\end{verbatim}

Error example:

\begin{verbatim}
{
  "timestamp": "2026-06-03T18:37:59Z",
  "model": "nvidia/nemotron-nano-12b-v2-vl:free",
  "target": "Bank Neo Sustainable Report/batch_27",
  "target_pages": "page_0052.md, page_0053.md",
  "prompt": "tone_few_shot_indonesian.md",
  "ok": false,
  "records": [],
  "error": "OpenRouter returned HTTP 401: {\"error\":{\"message..",
  "error_type": "unknown",
  "raw_output": "",
  "background_job_id": "llm_processing_bg_20260603T183739Z_855dc4"
}
\end{verbatim}

\subsection{ClimateBERT Result JSON}\label{climatebert-result-json}

Path example:

\begin{itemize}
\tightlist
\item
  \texttt{results/climatebert\_results.json}
\end{itemize}

Purpose:

\begin{itemize}
\tightlist
\item
  stores ClimateBERT or remote-space comparison runs
\item
  keeps the original input text
\item
  preserves raw multi-model response text
\item
  supports later parsing into label-level comparison tables
\end{itemize}

Example:

\begin{verbatim}
{
  "timestamp": "2026-03-26T15:43:20.151168Z",
  "mode": "remote_space",
  "space_url": "https://darisdzakwanhoesien-climatebert..",
  "input_text": "Company's media and entertainment ..",
  "response_raw": "### ...\n### netzero-reduction\n•",
  "response_parsed": null
}
\end{verbatim}

\subsection{Other JSON Families Used by the Repository}\label{other-json-families-used-by-the-repository}

Additional JSON artifact families in the repository include the following. These are not all equally central to the thesis argument, but each one supports a specific part of the executable research workflow.

\begin{itemize}
\item
  \texttt{results/data/mapping.json} Purpose: stores label-normalization mappings used to standardize ESG and sentiment terminology. Example content includes mappings such as \texttt{"environmental":\ "E"}, \texttt{"governance":\ "G"}, and \texttt{"lingkungan":\ "E"}. This file acts as a normalization bridge between raw extracted labels and the thesis-facing schema.

\item
  \texttt{results/ground\_truth.json} Purpose: stores manually entered or review-oriented ground-truth style records. The sampled structure includes a timestamp, model name, source, input text, and a nested \texttt{result} object. In the current repository state, some entries still preserve model-side errors, which is useful because it shows that the reference-building layer keeps incomplete or failed cases visible rather than silently discarding them.
\item
  \texttt{results/predictions.json} Purpose: stores prediction outputs, especially from ClimateBERT-style or T1 comparison runs. The current file contains long source text segments together with model identifiers and prediction payloads. In practice, this file acts as a raw comparison layer before more compact summary tables are derived.

\item
  \texttt{results/t1\_results.json} Purpose: stores serialized T1-stage prediction results in JSON form. These records are similar to \texttt{predictions.json}, but they are part of the more formal T1 artifact family used by the benchmark side of the workflow. The sampled entries show model names, original text, and nested result objects.

\item
  \texttt{results/t2\_results.json} Purpose: stores T2-stage ABSA-style outputs. The sampled entries contain both a \texttt{rule\_based} block and a \texttt{hybrid} block. This is important because it shows that T2 is not one monolithic classifier. It compares rule-based aspect/polarity/tone output against a hybrid contextual model that also records ontology alignment and summary metrics.
\end{itemize}

Together, these JSON families support benchmark generation, ontology mapping, workflow documentation, summarization inputs, transfer-learning dataset reporting, and interactive tooling. They form an important part of the repository's reproducibility and audit layer even when they are not all directly cited in the main thesis chapters.
 
\end{document}
